% PS6--9 source ledger: research/ps06-09-coverage.json
\originalchapter{作业六：周期响应、阶跃与卷积}
本章保留原作业的讲次题号。$u$ 表示单位阶跃，$\delta$ 表示单位冲激；响应均取因果响应，卷积取 $\int_0^t f(t-\tau)g(\tau)\,d\tau$。第二部分以官方解为依据逐问核对，补充必要推导；明确指出的补解及修正不冒充原答案。
\originalsection{第一部分}
\begin{exercise}\label{cw:ps6-I-22}原题 I.22（3 月 31 日）：完成 Notes 7C-1、7C-2。\end{exercise}
\begin{solution}
\textit{编者补解及题库回指核对（AI 辅助）；教材题面缺项另标。}\par

完整题面和解答见习题 \ref{cw:bank-7C-1} 的 (a)、(b)、(c) 及习题 \ref{cw:bank-7C-2}。前者依次检查三种强迫的谐波与固有频率是否重合；后者求周期特解。此处保留原指定及子问映射，不重复排印同一题。
\end{solution}
原题 I.23、I.25 均写明无题。
\begin{exercise}\label{cw:ps6-I-24}原题 I.24（4 月 5 日）：分别求 (a) $D+kI$、(b) $D^2+\omega_n^2I$ 的单位冲激响应与单位阶跃响应。\end{exercise}
\begin{solution}
\textit{原题附答译审与补证（AI 辅助）。}\par

(a) 对 $t>0$，齐次方程为 $w'+kw=0$；跨过零点积分得 $w(0+)=1$，故
\[w=u(t)e^{-kt},\qquad v=\int_0^tw(\tau)\,d\tau=\begin{cases}u(t)(1-e^{-kt})/k,&k\ne0,\\u(t)t,&k=0.\end{cases}\]
(b) 在 $\omega_n>0$ 时，冲激不使位移跳跃，但使速度跳跃 $1$，所以 $w(0+)=0,w'(0+)=1$，得到
\[w=u(t)\frac{\sin\omega_nt}{\omega_n},\qquad v=u(t)\frac{1-\cos\omega_nt}{\omega_n^2}.\]
直接求导得到 $v'=w$；$v(0+)=v'(0+)=0$，且 $v''+\omega_n^2v=1$。若 $\omega_n=0$，连续极限分别为 $u(t)t$、$u(t)t^2/2$。
\end{solution}
\originalsection{第二部分}
\begin{exercise}\label{cw:ps6-II-22}原题 II.22（周期解）：$g$ 的周期为 $2\pi$，在 $-\pi/2\le t\le\pi/2$ 时 $g(t)=t$，在 $\pi/2\le t\le3\pi/2$ 时 $g(t)=\pi-t$（即前一作业 II.21(f) 的函数）。
(a) 求 $x''+\omega_n^2x=g(t)$ 的一个周期解，若存在。
(b) 哪些正 $\omega_n$ 使周期解不存在？
(c) 将 (b) 最小值记作 $\omega_r$，当 $\omega_n$ 略小于、略大于 $\omega_r$ 时，解大致如何？
(d) 哪些 $\omega_n$ 使周期解不唯一？
(e) 对 (d) 的值，是否所有解都周期？\end{exercise}
\begin{solution}
\textit{官方解译审与补证（AI 辅助）；其中另标编者补解或原文修正。}\par

由奇性与分段积分，$g$ 只有奇数正弦谐波，且
\[g(t)=\frac4\pi\sum_{j=0}^{\infty}\frac{(-1)^j}{(2j+1)^2}\sin((2j+1)t).\]
(a) 若 $\omega_n$ 不是正奇整数，可逐谐波除以 $\omega_n^2-(2j+1)^2$：
\[x_p(t)=\frac4\pi\sum_{j=0}^{\infty}\frac{(-1)^j\sin((2j+1)t)}{(2j+1)^2[\omega_n^2-(2j+1)^2]}.\]
系数为 $O(j^{-4})$，两次导数的系数为 $O(j^{-2})$，因此上述级数及两次求导均一致收敛，所得为实际 $2\pi$ 周期解。
(b) 若 $\omega_n=N$ 为正奇整数，强迫的 $\sin Nt$ 系数非零。相应特解有非零的 $t\cos Nt$ 项，齐次解只有有界正弦、余弦，无法消除此无界项，故无周期解。
(c) $\omega_r=1$。靠近 $1$ 时第一项支配：
\[x_p(t)\approx\frac{4\sin t}{\pi(\omega_n^2-1)}.\]
左侧系数负，近似反相且振幅很大；右侧系数正，近似同相且振幅很大。其余谐波在邻域内一致有界。
(d) 一般解为 $x_p+A\cos\omega_nt+B\sin\omega_nt$。若非零齐次项与 $x_p$ 共存，则其频率与强迫的基本频率须可公度。严格地说，若和的周期为 $T$，则其方程右端也以 $T$ 为周期，三角波的周期必须是 $2\pi$ 的整数倍；再由齐次项周期得到 $\omega_nT\in2\pi\mathbb Z$。所以恰在 $\omega_n\in\mathbb Q_{>0}$ 且不是正奇整数时，有多个周期解。
(e) 是。写 $\omega_n=p/q$（互素正整数），所有解均以 $2\pi q$ 为周期。若 $\omega_n$ 为正偶整数，$2\pi$ 已是所有解的周期；若 $\omega_n$ 无理，只有 $x_p$ 这一周期解。
\end{solution}
\begin{exercise}\label{cw:ps6-II-23}原题 II.23（阶跃与冲激）：对下列每个 $f$，(i) 画函数图；(ii) 画广义导数图；(iii) 用 $u(t-a)$、$\delta(t-a)$ 及普通函数写出 $f$、$f'$（少数点可不定义）。
(a) $t<0$ 时 $f=0$，$t>0$ 时 $f=-t$。
(b) $t<0$ 时 $f=0$，$t>0$ 时 $f=1-t$。
(c) $t<0$ 时 $f=0$，$0<t<1$ 时 $f=2t-1$，$t>1$ 时 $f=0$。
(d) $t<0$ 时 $f=0$，$t>0$ 时 $f=t-\lfloor t\rfloor$，$\lfloor t\rfloor$ 是不超过 $t$ 的最大整数。\end{exercise}
\begin{solution}
\textit{官方解译审与补证（AI 辅助）；其中另标编者补解或原文修正。}\par

广义导数由各光滑区间的斜率与各跳点的跳跃量组成：$f(a+)-f(a-)$ 乘 $\delta(t-a)$。于是
\begin{align*}
\text{(a)}\quad&f=-tu(t),& f'&=-u(t),\\
\text{(b)}\quad&f=(1-t)u(t),& f'&=-u(t)+\delta(t),\\
\text{(c)}\quad&f=(2t-1)[u(t)-u(t-1)],&f'&=2[u(t)-u(t-1)]-\delta(t)-\delta(t-1),\\
\text{(d)}\quad&f=tu(t)-\sum_{n=1}^{\infty}u(t-n),&f'&=u(t)-\sum_{n=1}^{\infty}\delta(t-n).
\end{align*}
(d) 的和在任何有界区间都只有有限非零项；$0$ 处连续，不产生冲激，每个正整数处从 $1$ 跳到 $0$，故冲激权重为 $-1$。下图将冲激画为带权箭头，箭头不是普通函数值。
\begin{center}
\begin{tikzpicture}[x=1.1cm,y=.65cm,>=Stealth,font=\small]
\foreach \j/\title in {0/{(a)},1/{(b)},2/{(c)},3/{(d)}}{
\begin{scope}[shift={(3.1*\j,0)}]
\draw[->](-.3,0)--(2.4,0) node[right]{$t$};\draw[->](0,-1.3)--(0,1.5);\node at (1,1.8){\title：$f$};
\end{scope}}
\draw[main,thick](0,0)--(1.2,-1.2);
\draw[main,thick](3.1,1)--(5.1,-1);
\draw[main,thick](6.2,-1)--(7.2,1);\draw[main,thick](7.2,0)--(8.4,0);
\foreach \n in {0,1}{\draw[main,thick](9.3+\n,0)--(10.3+\n,1);}\node at (11.5,.5){$\cdots$};
\foreach \j/\title in {0/{(a)},1/{(b)},2/{(c)},3/{(d)}}{
\begin{scope}[shift={(3.1*\j,-4)}]\draw[->](-.3,0)--(2.4,0) node[right]{$t$};\draw[->](0,-1.3)--(0,1.5);\node at (1,1.8){\title：$f'$};\end{scope}}
\draw[main,thick](0,-5)--(2.2,-5);
\draw[main,thick](3.1,-5)--(5.3,-5);\draw[->,thick](3.1,-4)--(3.1,-3) node[above right]{$+1$};
\draw[main,thick](6.2,-2.7)--(7.2,-2.7);\draw[main,thick](7.2,-4)--(8.4,-4);\foreach \x in {6.2,7.2}{\draw[->,thick](\x,-4)--(\x,-5) node[below]{$-1$};}\node at (6.7,-2.4){$2$};
\draw[main,thick](9.3,-3.3)--(11.5,-3.3);\foreach \x in {10.3,11.3}{\draw[->,thick](\x,-4)--(\x,-5) node[below]{$-1$};}\node at (9.7,-3){$1$};
\end{tikzpicture}
\end{center}
\end{solution}
\begin{exercise}\label{cw:ps6-II-24}原题 II.24：
(a) 求 $2D^2+4D+4I$ 的单位冲激响应 $w$。
(b) 求其单位阶跃响应 $v$。
(c) 验证 $v'=w$。
(d) 分别求以 (i) $2u(t)$、(ii) $u(t)t$、(iii) $u(t)t^2$ 为单位冲激响应的 LTI 微分算子 $p(D)$。\end{exercise}
\begin{solution}
\textit{官方解译审与补证（AI 辅助）；其中另标编者补解或原文修正。}\par

(a) 特征根为 $-1\pm i$，故 $t>0$ 时 $w=e^{-t}(A\cos t+B\sin t)$。冲激要求 $w(0+)=0$、$2w'(0+)=1$，于是
\[w=\tfrac12u(t)e^{-t}\sin t.\]
(b) 常数特解为 $1/4$；将齐次项加上后，以静止条件 $v(0+)=v'(0+)=0$ 确定系数：
\[v=\tfrac14u(t)[1-e^{-t}(\cos t+\sin t)].\]
(c) 因括号在零点为零，求广义导数不产生冲激，且
\[v'=\tfrac14u(t)e^{-t}[(\cos t+\sin t)-(\cos t-\sin t)]=w.\]
(d) 直接应用广义导数：(i) $D(2u)=2\delta$，故 $p(D)=D/2$；(ii) $D(ut)=u$、$D^2(ut)=\delta$，故 $p(D)=D^2$；(iii) $D(ut^2)=2ut$、$D^2(ut^2)=2u$、$D^3(ut^2)=2\delta$，故 $p(D)=D^3/2$。这些算子的低阶系数必须为零，因为在 $t>0$ 上作用所得多项式恒为零；更高阶非零项会产生额外的冲激导数，亦不允许。
\end{solution}
\begin{exercise}\label{cw:ps6-II-25}原题 II.25（卷积）：
(a) 令 $q(t)=\cos\omega t$，用 I.24 的 $D+kI$ 冲激响应求 $w*q$，验证它是 $x'+kx=q$ 的静止初值解。
(b) 令 $q(t)=1$，用 I.24 的 $D^2+\omega_n^2I$ 冲激响应求 $w*q$，验证它是 $x''+\omega_n^2x=q$ 的静止初值解。
(c) 计算 $t^2*t$、$t*t^2$，比较。
(d) 计算 $(t*t)*t$、$t*(t*t)$，比较。\end{exercise}
\begin{solution}
\textit{官方解译审与补证（AI 辅助）；其中另标编者补解或原文修正。}\par

(a) 当 $k^2+\omega^2>0$ 时，
\begin{align*}
x(t)&=e^{-kt}\operatorname{Re}\int_0^t e^{(k+i\omega)\tau}\,d\tau\\
&=\frac{k\cos\omega t+\omega\sin\omega t-ke^{-kt}}{k^2+\omega^2}.
\end{align*}
其导数为 $[-k\omega\sin\omega t+\omega^2\cos\omega t+k^2e^{-kt}]/(k^2+\omega^2)$，加 $kx$ 恰为 $\cos\omega t$，而 $x(0)=0$。$k=\omega=0$ 时直接积分得 $x=t$。
(b) $x=\omega_n^{-1}\int_0^t\sin\omega_n(t-\tau)\,d\tau=(1-\cos\omega_nt)/\omega_n^2$。于是 $x'=\sin\omega_nt/\omega_n$、$x''=\cos\omega_nt$，方程与两个零初值均成立。
(c) 展开被积式：
\[t^2*t=\int_0^t(t-\tau)^2\tau\,d\tau=(\tfrac12-\tfrac23+\tfrac14)t^4=\frac{t^4}{12},\]
\[t*t^2=\int_0^t(t-\tau)\tau^2\,d\tau=(\tfrac13-\tfrac14)t^4=\frac{t^4}{12}.\]
(d) $t*t=\int_0^t(t-\tau)\tau\,d\tau=t^3/6$。因此
\[(t*t)*t=\tfrac16\int_0^t(t-\tau)^3\tau\,d\tau=\frac{t^5}{120},\quad t*(t*t)=\tfrac16\int_0^t(t-\tau)\tau^3\,d\tau=\frac{t^5}{120}.\]
这些计算分别展示交换律与结合律，所有输入在此均按 $t\ge0$ 的因果函数理解。
\end{solution}
\originalchapter{作业七：拉普拉斯变换}
\originalsection{第一部分}
\begin{exercise}\label{cw:ps7-I-26}原题 I.26（4 月 9 日）：令 $z$ 为复数。从定义求 $\mathcal L[e^{zt}]$，并求积分的收敛域。\end{exercise}
\begin{solution}
\textit{编者补解及题库回指核对（AI 辅助）；教材题面缺项另标。}\par

\[\int_0^R e^{-(s-z)t}\,dt=\frac{1-e^{-(s-z)R}}{s-z}\quad(s\ne z).\]
极限存在恰需 $\operatorname{Re}(s-z)>0$，此时 $F(s)=1/(s-z)$。若实部小于零，边界项增长；若实部等于零而 $s\ne z$，边界项振荡无极限；$s=z$ 时积分为 $R$。因此收敛域为 $\operatorname{Re}s>\operatorname{Re}z$。
\end{solution}
\begin{exercise}\label{cw:ps7-I-27}原题 I.27（4 月 12 日）：完成 Notes 3A-9（(a) 先按建议用 $s$ 平移规则，再将 $e^{-t}\sin3t$ 写成复指数的线性组合重做）、Notes 3B-1；EP 第 4.3 节第 1、5 题。\end{exercise}
\begin{solution}
\textit{编者补解及题库回指核对（AI 辅助）；教材题面缺项另标。}\par

Notes 3A-9 的 (a)、(b) 见习题 \ref{cw:bank-3A-9}，Notes 3B-1 的 (a)--(e) 见习题 \ref{cw:bank-3B-1}。本作业额外要求的重做为
\[e^{-t}\sin3t=\frac{e^{(-1+3i)t}-e^{(-1-3i)t}}{2i},\]
\[\mathcal L[e^{-t}\sin3t]=\frac1{2i}\left(\frac1{s+1-3i}-\frac1{s+1+3i}\right)=\frac3{(s+1)^2+9},\quad\operatorname{Re}s>-1.\]
EP 4.3 第 1、5 题仅有教材定位，未取得可核实的公开题面（\texttt{external-textbook-statement-unavailable}），故此处不猜写题面或解答。
\end{solution}
\begin{exercise}\label{cw:ps7-I-28}原题 I.28（4 月 14 日）：完成 EP 第 4.3 节第 7、8、12、13、28、36 题。\end{exercise}
\begin{solution}
\textit{编者补解及题库回指核对（AI 辅助）；教材题面缺项另标。}\par

以上六项保留完整指定；公开作业未包含它们的题面，当前公开题源亦未取得相应教材正文，状态均为 \texttt{external-textbook-statement-unavailable}。不能用题号推测或借另一版教材同号题代替。
\end{solution}
\originalsection{第二部分}
\begin{exercise}\label{cw:ps7-II-26}原题 II.26：
(a) $F(s)=\mathcal L[f](s)$，$a>0$。用积分定义及换元求 $g(t)=f(at)$ 的变换，并用 $f(t)=t^n$ 的两边变换验证。
(b) 用二重积分及换元 $x=t-\tau,y=\tau$ 证明 $h=f*g$ 时 $H=FG$；从 $F=\int_0^\infty f(x)e^{-sx}\,dx$、$G=\int_0^\infty g(y)e^{-sy}\,dy$ 出发。
(c) 原题写作「$0<t<1$ 时 $f(t)=1$，$t>0$ 时 $f(t)=0$」，要求从定义求变换及收敛域。\end{exercise}
\begin{solution}
\textit{官方解译审与补证（AI 辅助）；其中另标编者补解或原文修正。}\par

(c) 原文两个条件在 $(0,1)$ 相互矛盾；官方解将后一条件按 $t>1$ 处理。以下明确采用该修正，即单位矩形脉冲，保留这一差异记录。
(a) 令 $v=at$，则
\[G(s)=\int_0^\infty f(at)e^{-st}\,dt=\frac1a\int_0^\infty f(v)e^{-(s/a)v}\,dv=\frac1aF(s/a).\]
若 $F$ 在 $\operatorname{Re}s>\sigma$ 收敛，$G$ 在 $\operatorname{Re}s>a\sigma$ 收敛。$n$ 为非负整数时，$F=n!/s^{n+1}$、$G=a^nn!/s^{n+1}$，而 $a^{-1}F(s/a)=a^nn!/s^{n+1}$，一致。
(b) 在 $\operatorname{Re}s$ 足够大、两积分绝对收敛时，Fubini 定理允许写
\[FG=\int_0^\infty\int_0^\infty f(x)g(y)e^{-s(x+y)}\,dx\,dy.\]
新区域为 $0\le\tau\le t<\infty$，Jacobi 行列式 $\det\begin{pmatrix}1&-1\\0&1\end{pmatrix}=1$。故
\[FG=\int_0^\infty e^{-st}\left[\int_0^t f(t-\tau)g(\tau)\,d\tau\right]dt=H(s).\]
绝对收敛的假设确保这里的乘积、换序与换元合法。
(c) 修正后的变换为
\[F(s)=\int_0^1 e^{-st}\,dt=\begin{cases}(1-e^{-s})/s,&s\ne0,\\1,&s=0.\end{cases}\]
这是有限区间积分，对全部复 $s$ 收敛；$s=0$ 是表达式的可去奇点，并非极点。
\end{solution}
\begin{exercise}\label{cw:ps7-II-27}原题 II.27：$a,b$ 为实数，$a\ne0$。
(a) 对 $aD+bI$，通过拉普拉斯变换到频域求解再变回，求单位冲激响应及阶跃响应。宜使用因果广义导数规则 $\mathcal L[x']=sX$。
(b) 用三种方法求 $ax'+bx=t$ 的静止初值解：(i) 待定系数加适当瞬态；(ii) 用 (a) 的 $w$ 求 $w*t$；(iii) 拉普拉斯变换。\end{exercise}
\begin{solution}
\textit{官方解译审与补证（AI 辅助）；其中另标编者补解或原文修正。}\par

(a) $(as+b)W=1$，故
\[W=\frac1{as+b},\quad w=\frac1a u(t)e^{-bt/a},\quad V=\frac1{s(as+b)}.\]
当 $b\ne0$ 时，$V=b^{-1}(s^{-1}-(s+b/a)^{-1})$，所以 $v=u(t)(1-e^{-bt/a})/b$；当 $b=0$ 时 $v=u(t)t/a$。
(b)(i) 设 $x_p=ct+d$，得到 $bc=1$、$ac+bd=0$。$b\ne0$ 时 $x_p=t/b-a/b^2$，加 $Ce^{-bt/a}$ 并令 $x(0)=0$ 得
\[x=u(t)\left[\frac tb-\frac a{b^2}(1-e^{-bt/a})\right].\]
$b=0$ 时直接积分 $x'=t/a$，得 $x=u(t)t^2/(2a)$。
(ii) $b\ne0$ 时，分部积分给出
\begin{align*}
(w*t)(t)&=\frac1a e^{-bt/a}\int_0^t\tau e^{b\tau/a}\,d\tau\\
&=\frac1a e^{-bt/a}\left[\frac{a\tau}{b}e^{b\tau/a}-\frac{a^2}{b^2}e^{b\tau/a}\right]_0^t\\
&=\frac tb-\frac a{b^2}(1-e^{-bt/a}).
\end{align*}
$b=0$ 时积分为 $a^{-1}\int_0^t\tau\,d\tau=t^2/(2a)$。
(iii) $X=1/[s^2(as+b)]$。$b\ne0$ 时分解为
\[X=-\frac a{b^2s}+\frac1{bs^2}+\frac a{b^2(s+b/a)},\]
反变换得同一公式；$b=0$ 时 $X=1/(as^3)$。在 $t>0$，$x'=(1-e^{-bt/a})/b$，代入 $ax'+bx=t$，同时 $x(0+)=0$；三种构造的解确实一致。
\end{solution}
\begin{exercise}\label{cw:ps7-II-28}原题 II.28：用拉普拉斯变换求以下算子的单位冲激响应：(a) $3D^2+6D+6I$；(b) $D^4+I$。\end{exercise}
\begin{solution}
\textit{官方解译审与补证（AI 辅助）；其中另标编者补解或原文修正。}\par

(a) $W=[3((s+1)^2+1)]^{-1}$，故 $w=u(t)e^{-t}\sin t/3$。
(b) 原题确为 $D^4+I$。官方解只列出四个根及部分分式框架，并说明原本想问 $D^4-I$；下面是对实际题面的编者补解。令 $\alpha=1/\sqrt2$，则
\[s^4+1=(s^2-\sqrt2s+1)(s^2+\sqrt2s+1),\]
\[\frac1{s^4+1}=\frac{-s/(2\sqrt2)+1/2}{s^2-\sqrt2s+1}+\frac{s/(2\sqrt2)+1/2}{s^2+\sqrt2s+1}.\]
将两分母写成 $(s\mp\alpha)^2+\alpha^2$，反变换并合并指数项，得
\[w(t)=\frac{u(t)}{\sqrt2}\left[\cosh(\alpha t)\sin(\alpha t)-\sinh(\alpha t)\cos(\alpha t)\right].\]
四个指数根为 $\pm\alpha\pm i\alpha$，每项在 $t>0$ 满足 $w^{(4)}+w=0$。更直接地，括号的展开从 $(2/3)(\alpha t)^3$ 开始，故 $w=t^3/6+O(t^7)$，于是 $w(0+)=w'(0+)=w''(0+)=0$、$w'''(0+)=1$。因果延拓的四阶导数恰产生单位冲激，验证了所求响应。若改问官方所述的 $D^4-I$，答案才是 $u(t)(\sinh t-\sin t)/2$；二者不可混用。
\end{solution}
\originalchapter{作业八：极点与线性系统}
原作业分两半发布，以下合并但保留讲次 29、32、33；讲次 30 无题、31 为第三次考试，不补造题目。第一半将 L32 日期写成 4 月 28 日，第二半实际题目写 4 月 26 日，以下依第二半。
\originalsection{第一部分}
\begin{exercise}\label{cw:ps8-I-29}原题 I.29：
(a) 用 $t$ 平移规则求 $f(t)=[u(t)-u(t-2\pi)]\sin t$ 的拉普拉斯变换。
(b) 画下列函数变换的极点图：(i) $f=1$；(ii) $f=e^{-t}+3e^{-3t}$；(iii) $f=\cos2t+e^{-t}\sin t$。\end{exercise}
\begin{solution}
\textit{原题附答译审与补证（AI 辅助）。}\par

(a) 令 $g=u(t)\sin t$，则 $g(t-2\pi)=u(t-2\pi)\sin t$，所以
\[F(s)=\frac{1-e^{-2\pi s}}{s^2+1}.\]
虽然分母在 $\pm i$ 为零，分子亦为零且抵消；$f$ 有紧支集，$F$ 在整个复平面解析，故这两个点不是极点。
(b) 三个变换分别为
\[\frac1s,\quad \frac1{s+1}+\frac3{s+3},\quad\frac{s}{s^2+4}+\frac1{(s+1)^2+1}.\]
无抵消，其极点依次为 $\{0\}$、$\{-1,-3\}$、$\{\pm2i,-1\pm i\}$。极点图如下，横轴为实部、纵轴为虚部，叉号为极点。
\begin{center}\begin{tikzpicture}[x=.8cm,y=.65cm,font=\small]
\foreach \a/\ttl in {0/{(i)},5/{(ii)},11/{(iii)}}{\begin{scope}[shift={(\a,0)}]\draw[->](-3.3,0)--(1,0) node[right]{$\operatorname{Re}s$};\draw[->](0,-2.3)--(0,2.5) node[above]{$\operatorname{Im}s$};\node at (-1,3){\ttl};\end{scope}}
\node[main] at (0,0){$\times$};\node[main] at (2,0){$\times$};\node[main] at (4,0){$\times$};\node[below] at (2,0){$-3$};\node[below] at (4,0){$-1$};
\foreach \x/\y in {11/2,11/-2,10/1,10/-1}{\node[main] at (\x,\y){$\times$};}\node[right] at (11,2){$2i$};\node[right] at (11,-2){$-2i$};\node[above left] at (10,1){$-1+i$};\node[below left] at (10,-1){$-1-i$};
\end{tikzpicture}\end{center}
\end{solution}
\begin{exercise}\label{cw:ps8-I-32}原题 I.32：
(a) 计算五个矩阵乘积
\[\begin{pmatrix}1&2\end{pmatrix}\binom xy,\quad\binom12\begin{pmatrix}x&y\end{pmatrix},\quad\begin{pmatrix}a&b\\c&d\end{pmatrix}\binom xy,\quad\begin{pmatrix}1&2\end{pmatrix}\begin{pmatrix}a&b\\c&d\end{pmatrix},\quad\begin{pmatrix}a&b\\c&d\end{pmatrix}\begin{pmatrix}x&u\\y&v\end{pmatrix}.\]
(b) 完成 Notes 4A-2、4B-1、4B-6。\end{exercise}
\begin{solution}
\textit{编者补解及题库回指核对（AI 辅助）；教材题面缺项另标。}\par

(a) 每一项用左矩阵的一行与右矩阵的一列作内积，依次得
\[x+2y,\quad \begin{pmatrix}x&y\\2x&2y\end{pmatrix},\quad\binom{ax+by}{cx+dy},\quad\begin{pmatrix}a+2c&b+2d\end{pmatrix},\quad\begin{pmatrix}ax+by&au+bv\\cx+dy&cu+dv\end{pmatrix}.\]
(b) 见习题 \ref{cw:bank-4A-2}、\ref{cw:bank-4B-1} 的 (a)、(b)、\ref{cw:bank-4B-6} 的 (a)、(b)，后一题两种消元方法须比较。
\end{solution}
\begin{exercise}\label{cw:ps8-I-33}原题 I.33：完成 Notes 4C-1；可参考 EP 第 344 页的 $3\times3$ 行列式计算说明。\end{exercise}
\begin{solution}
\textit{编者补解及题库回指核对（AI 辅助）；教材题面缺项另标。}\par
完整的三个矩阵题及解答见习题 \ref{cw:bank-4C-1} 的 (a)、(b)、(c)。EP 页码在此只作为计算方法的参考指定，不作为另一个未给题面的作业题。\end{solution}
\originalsection{第二部分}
\begin{exercise}\label{cw:ps8-II-29}原题 II.29（极点）：
(a) 对四个极点图 (1) $\{1,i,-i\}$；(2) $\{-1+4i,-1-4i\}$；(3) $\{-1\}$；(4) 空图，分别 (i) 描述具有该极点图的所有 $f$ 的共同特征；(ii) 写出两个 $f,F$ 的例子。
(b) 埃塞俄比亚的一次考古发掘发现一套机械系统，为避免拆开，研究者施加单位冲激，测得 $w(t)=u(t)e^{-t/2}\sin(3t/2)$。
(i) 求系统的特征多项式与传递函数 $W(s)$。
(ii) 画极点图。
(iii) 运送卡车的振动会激起系统振动，应避开哪个使响应振幅最大的圆频率？原题此处再次标成 (ii)，此册按第三小问记。
(iv) 在 Mathlet「Amplitude Response and Pole Diagram」中设定系统参数，检查峰值并转动三维图，解释黄色、绿色盒状图，黄色盒底曲线，红箭头及黄菱形。第二半纠正：Mathlet 的 $b,k$ 应输入 (i) 的 $b/m,k/m$。\end{exercise}
\begin{solution}
\textit{官方解译审与补证（AI 辅助）；其中另标编者补解或原文修正。}\par

(a) 需要先说明推断范围。若 $F$ 为严格真有理函数、图中各极点均为单极点、$f$ 为实函数，则部分分式反演给出以下共同结构。仅列极点位置、不列阶数或允许非有理变换时，不能声称这些式子描述「所有函数」：重极点引入 $t$ 的多项式因子，整函数项还可加入无新极点的时域成分。以下结构与官方解采用的简单指数模式相同，同时说明其适用范围。
(1) $f=Ae^t+B\cos t+C\sin t$，$A\ne0$ 且 $(B,C)\ne(0,0)$。长期增长由 $e^t$ 支配，另有圆频率 $1$ 的恒振幅振荡。两个例子为
\[f_1=e^t+\sin t,\quad F_1=\frac1{s-1}+\frac1{s^2+1};\qquad f_2=2e^t+\cos t,\quad F_2=\frac2{s-1}+\frac{s}{s^2+1}.\]
(2) $f=e^{-t}(B\cos4t+C\sin4t)$，$(B,C)\ne(0,0)$，是包络按 $e^{-t}$ 衰减的圆频率 $4$ 振荡。例子：
\[f_1=e^{-t}\sin4t,\ F_1=\frac4{(s+1)^2+16};\qquad f_2=e^{-t}\cos4t,\ F_2=\frac{s+1}{(s+1)^2+16}.\]
(3) $f=Ae^{-t}$、$A\ne0$，无振荡并指数衰减；可取 $A=1,2$，相应 $F=1/(s+1),2/(s+1)$。
(4) 空图只表明变换无极点，不能仅据此断定每种 $f$ 的长期衰减速度。若限于上述严格真有理类，只能有零函数；允许一般因果函数或广义函数时，可取
\[f_1=\delta(t-1),\quad F_1=e^{-s};\qquad f_2=u(t)-u(t-1),\quad F_2=\frac{1-e^{-s}}s,\quad F_2(0)=1.\]
两者均在有限时间后为零，但这只是这些例子的性质。官方解将空极点图概括成「比任意指数衰减更快」，此结论需要额外函数类假设，不从极点位置单独推出。
(b)(i) 变换给出
\[W(s)=\frac{3/2}{(s+1/2)^2+9/4}=\frac{3/2}{s^2+s+5/2},\quad p(s)=\frac23s^2+\frac23s+\frac53.\]
因此机械参数为 $m=2/3,b=2/3,k=5/3$。也可从 $w'(0+)=3/2=1/m$ 确定首项系数，不能只写首一特征多项式而丢失传递函数的增益。
(ii) 极点为 $-1/2\pm3i/2$。
(iii) 对单位幅度谐波，增益
\[|W(i\omega)|=\frac{3/2}{\sqrt{(5/2-\omega^2)^2+\omega^2}}.\]
令 $z=\omega^2\ge0$，分母平方为 $(5/2-z)^2+z=(z-2)^2+9/4$，故唯一正峰值频率为 $\omega_r=\sqrt2$，最大增益为 $1$。
(iv) 以下依据函数和官方图例作解析解释，未声称实际操作了 Mathlet。输入其归一化参数 $b/m=1,k/m=5/2$；其首一系统增益与实际 $W$ 相差常数 $2/3$，峰值频率一致。黄色盒状图处于虚轴上方的竖直平面，按官方图例，其底边界曲线是振幅响应 $(i\omega,|W(i\omega)|)$。绿色盒状图处于实轴上方的平面，顶部是三维曲面 $|W(s)|$ 在实轴上的截线、底部是实轴。黄色盒的底曲线将输入频率与相应增益配对，并不是虚轴本身，红箭头位于 $W$ 的极点上方，黄菱形在 $(\pm i\omega,|W(i\omega)|)$ 处标记所选输入频率与相应增益。
\begin{center}\begin{tikzpicture}
\begin{axis}[width=6cm,height=4.8cm,xmin=-1.6,xmax=.5,ymin=-2,ymax=2,axis lines=middle,xlabel={$\operatorname{Re}s$},ylabel={$\operatorname{Im}s$},title={极点图},xtick={-.5},ytick={-1.5,1.5}]
\addplot[only marks,mark=x,mark size=4pt,main] coordinates{(-.5,-1.5)(-.5,1.5)};
\end{axis}
\begin{axis}[at={(6.2cm,0)},width=6cm,height=4.8cm,xmin=0,xmax=4,ymin=0,ymax=1.15,axis lines=left,xlabel={$\omega$},ylabel={增益},title={实际系统振幅响应},xtick={0,1.4142,3},xticklabels={$0$,$\sqrt2$,$3$}]
\addplot[main,thick,domain=0:4,samples=120]{1.5/sqrt((2.5-x*x)^2+x*x)};\addplot[only marks] coordinates{(1.41421356,1)};
\end{axis}\end{tikzpicture}\end{center}
\end{solution}
\begin{exercise}\label{cw:ps8-II-32}原题 II.32（线性系统与矩阵）：研究 $x''+4x'+3x=0$、$x''+x'+\tfrac52x=0$。
(a) 每个方程求两个线性无关实解 $x_1,x_2$（取指数或 $e^{rt}\cos\omega t,e^{rt}\sin\omega t$），写一般实解、阻尼类型，并算 $x_1',x_2'$。
(b) 写两个伴随矩阵。
(c) 在 Mathlet「Linear Phase Portraits: Matrix Entry」选择伴随矩阵，将 $c,d$ 设为第一方程的对应值；可在左上 $(d,-c)$ 平面调整，点击相平面增添轨道，再清除并恢复初始轨道。画相平面，标时间箭头，指出 (a) 的两个基本解所对应的轨道。
(d) 若图示最大 $x=2$，一条钩形轨道在 $t=0$ 经过 $x$ 轴的 $x=1$ 点，求该解并画原方程的 $x(t)$。
(e) 写同一轨道上的一般解；可设在 $t=a$ 时 $x(a)=1$。
(f) 将 $c,d$ 改成第二方程值，画带时间箭头的相图。同一尺度下，求 $t=0$ 经过 $(0,1)$ 的解；何时再次过 $y$ 轴？粗画 $x(t),y(t)$。
(g) 解释这些伴随矩阵的轨道过 $x$ 轴时切向量为何竖直。\end{exercise}
\begin{solution}
\textit{官方解译审与补证（AI 辅助）；其中另标编者补解或原文修正。}\par

(a) 第一式特征多项式 $(r+1)(r+3)$，可取
\[x_1=e^{-t},\ x_2=e^{-3t},\quad x_1'=-e^{-t},\ x_2'=-3e^{-3t},\quad x=C_1e^{-t}+C_2e^{-3t};\]
是过阻尼。第二式根为 $-1/2\pm3i/2$，可取
\[x_1=e^{-t/2}\cos\frac{3t}2,\quad x_2=e^{-t/2}\sin\frac{3t}2,\quad x=C_1x_1+C_2x_2;\]
\[x_1'=e^{-t/2}\left(-\tfrac12\cos\tfrac{3t}2-\tfrac32\sin\tfrac{3t}2\right),\quad x_2'=e^{-t/2}\left(-\tfrac12\sin\tfrac{3t}2+\tfrac32\cos\tfrac{3t}2\right).\]
两根有非零虚部，是欠阻尼。两对解的初值向量均线性无关，故构成全部实解。
(b) 置 $y=x'$，伴随矩阵分别为
\[A_1=\begin{pmatrix}0&1\\-3&-4\end{pmatrix},\qquad A_2=\begin{pmatrix}0&1\\-5/2&-1\end{pmatrix}.\]
(c) 这里用解析轨道替代实际 Mathlet 操作的陈述。第一系统稳定结点，基本解向量为 $e^{-t}(1,-1)^T$、$e^{-3t}(1,-3)^T$；其轨道在 $y=-x$、$y=-3x$，箭头均指向原点。非特征线轨道也趋于原点，通常以慢衰减的 $y=-x$ 为到达切线。
(d) $x(0)=1,y(0)=0$ 给出 $C_1+C_2=1$、$-C_1-3C_2=0$，所以
\[x=\tfrac12(3e^{-t}-e^{-3t}),\qquad y=\tfrac32(-e^{-t}+e^{-3t}).\]
$x$ 在 $t=-\tfrac12\log3$ 过零，在 $t=0$ 达最大值 $1$，随后单调下降到 $0$；$x''=\tfrac32(e^{-t}-3e^{-3t})$ 在 $t=\tfrac12\log3$ 变号。过去 $t\to-\infty$ 时 $x\to-\infty$，故画出钩形转向。
(e) 自治方程的同轨道解是时间平移：
\[\binom{x(t)}{y(t)}=\frac12\binom{3e^{-(t-a)}-e^{-3(t-a)}}{-3e^{-(t-a)}+3e^{-3(t-a)}},\quad a\in\mathbb R.\]
(f) 第二系统为顺时针稳定螺旋：在正 $x$ 轴上速度向下。由 $x(0)=0,x'(0)=1$ 得
\[x=\frac23e^{-t/2}\sin\frac{3t}2,\qquad y=e^{-t/2}\left(\cos\frac{3t}2-\frac13\sin\frac{3t}2\right).\]
$x=0$ 恰在 $t=2n\pi/3$，$n\in\mathbb Z$；下一次为 $2\pi/3$。二者均为衰减振荡，$x(0)=0,x'(0)=1$，$y(0)=1,y'(0)=-1$。
(g) $x$ 轴上 $y=x'=0$，故 $\mathbf u'=(0,x'')^T=(0,cx)^T$，非原点处因本题 $c\ne0$ 而为非零竖直向量。原点是平衡点，不谈非零切向量。
\begin{center}\begin{tikzpicture}
\begin{axis}[width=6cm,height=5cm,xmin=-1.2,xmax=1.5,ymin=-2,ymax=2,axis lines=middle,xlabel={$x$},ylabel={$y=x'$},title={第一系统：稳定结点},ticks=none]
\addplot[dashed,domain=-1.2:1.2]{-x};\addplot[dashed,domain=-.66:.66]{-3*x};
\foreach \c/\d in {1/.2,1.5/-.5,-1/-.2,-1.5/.5}{\addplot[main,domain=-.25:3,samples=70,->]({\c*exp(-x)+\d*exp(-3*x)},{-\c*exp(-x)-3*\d*exp(-3*x)});}
\node at (axis cs:.8,-1.6){$y=-3x$};\node at (axis cs:1.1,-.8){$y=-x$};
\end{axis}
\begin{axis}[at={(6.2cm,0)},width=6cm,height=5cm,xmin=-1.2,xmax=1.2,ymin=-1.5,ymax=1.5,axis lines=middle,xlabel={$x$},ylabel={$y$},title={第二系统：顺时针螺旋},ticks=none]
\addplot[main,thick,domain=-.6:9,samples=160,->]({2/3*exp(-x/2)*sin(deg(1.5*x))},{exp(-x/2)*(cos(deg(1.5*x))-sin(deg(1.5*x))/3)});\addplot[only marks] coordinates{(0,1)};
\end{axis}\end{tikzpicture}
\begin{tikzpicture}
\begin{axis}[width=6cm,height=4.5cm,xmin=-.8,xmax=4,ymin=-1,ymax=1.2,axis lines=middle,xlabel={$t$},ylabel={$x$},title={(d) 位移曲线}]
\addplot[main,domain=-.8:4,samples=100]{(3*exp(-x)-exp(-3*x))/2};\addplot[only marks]coordinates{(0,1)};
\end{axis}
\begin{axis}[at={(6.2cm,0)},width=6cm,height=4.5cm,xmin=0,xmax=7,ymin=-.7,ymax=1.2,axis lines=middle,xlabel={$t$},title={(f) 时间曲线},legend style={font=\small}]
\addplot[main,domain=0:7,samples=120]{2/3*exp(-x/2)*sin(deg(1.5*x))};\addlegendentry{$x(t)$}
\addplot[structurecolor,dashed,domain=0:7,samples=120]{exp(-x/2)*(cos(deg(1.5*x))-sin(deg(1.5*x))/3)};\addlegendentry{$y(t)$}
\end{axis}\end{tikzpicture}\end{center}
\end{solution}
\begin{exercise}\label{cw:ps8-II-33}原题 II.33：
(a) 求 $x''+4x'+3x=0$ 的伴随矩阵特征值、特征向量，画特征线，各写一个沿线运动的解，与 II.32 特别是 (c) 比较。
(b) 对 $A=\begin{pmatrix}2&-4\\1&-3\end{pmatrix}$ 求特征值、特征向量，画特征线，各写出轨道处于该线的所有解。
(c) 在上述 Mathlet 中设该矩阵，画相平面，至少在每个特征线夹角区域各选一条轨道，并标时间箭头。\end{exercise}
\begin{solution}
\textit{官方解译审与补证（AI 辅助）；其中另标编者补解或原文修正。}\par

(a) $\det(\lambda I-A_1)=\lambda^2+4\lambda+3$，特征值 $-1,-3$。由第一行 $y=\lambda x$，特征向量可取 $(1,-1)^T,(1,-3)^T$；特征线及解正是 II.32(c) 的两条线和 $e^{-t}(1,-1)^T,e^{-3t}(1,-3)^T$。
(b) $\det(\lambda I-A)=\lambda^2+\lambda-2=(\lambda+2)(\lambda-1)$。$\lambda=-2$ 时向量可取 $(1,1)^T$，$\lambda=1$ 时可取 $(4,1)^T$。所有沿线解分别为 $Ce^{-2t}(1,1)^T$、$Ce^t(4,1)^T$，$C\in\mathbb R$。第一线向内，第二线向外。
(c) 全解为 $\mathbf u=C_1e^{-2t}(1,1)^T+C_2e^t(4,1)^T$。取 $(C_1,C_2)$ 的四组符号即得四个夹角的代表轨道；它们过去趋近稳定特征线方向、将来沿不稳定特征线逃离，形成鞍点。下图由此式构造，并未宣称运行 Mathlet。
\begin{center}\begin{tikzpicture}\begin{axis}[width=8cm,height=6cm,xmin=-4,xmax=4,ymin=-3,ymax=3,axis lines=middle,xlabel={$x$},ylabel={$y$},title={鞍点与四个夹角的轨道},ticks=none]
\addplot[dashed,domain=-3:3]{x};\addplot[dashed,domain=-4:4]{x/4};
\foreach \c/\d in {1/1,1/-1,-1/1,-1/-1}{\addplot[main,thick,domain=-.8:1.4,samples=90,->]({\c*exp(-2*x)+.4*\d*exp(x)},{\c*exp(-2*x)+.1*\d*exp(x)});}
\draw[->,thick](axis cs:2,2)--(axis cs:1,1);\draw[->,thick](axis cs:-2,-2)--(axis cs:-1,-1);\draw[->,thick](axis cs:1,.25)--(axis cs:2,.5);\draw[->,thick](axis cs:-1,-.25)--(axis cs:-2,-.5);
\node at (axis cs:1.7,2.5){$\lambda=-2$};\node at (axis cs:2.8,.4){$\lambda=1$};
\end{axis}\end{tikzpicture}\end{center}
\end{solution}
\originalchapter{作业九：相图分类与矩阵指数}
\originalsection{第一部分}
\begin{exercise}\label{cw:ps9-I-34}原题 I.34：完成 Notes 4D-2；EP 第 5.6 节第 1 题，可用「Linear Phase Portraits: Matrix Entry」辅助画相图。\end{exercise}
\begin{solution}
\textit{编者补解及题库回指核对（AI 辅助）；教材题面缺项另标。}\par
Notes 4D-2 的完整题面与解答见习题 \ref{cw:bank-4D-2}。EP 5.6 第 1 题只取得教材定位，未获可核实公开题面，标为 \texttt{external-textbook-statement-unavailable}，不以其他矩阵题代替。\end{solution}
\begin{exercise}\label{cw:ps9-I-35}原题 I.35：完成 Notes 5B-4，可用同一 Mathlet 检查。\end{exercise}
\begin{solution}
\textit{编者补解及题库回指核对（AI 辅助）；教材题面缺项另标。}\par
见习题 \ref{cw:bank-5B-4} 的 (a)--(e)，分别检查特征值及稳定性、实特征线或复根的旋转方向，再画带时间方向的轨道。\end{solution}
\begin{exercise}\label{cw:ps9-I-36}原题 I.36：完成 EP 第 5.7 节第 1、3 题。\end{exercise}
\begin{solution}
\textit{编者补解及题库回指核对（AI 辅助）；教材题面缺项另标。}\par
两项均只有教材题号指定，未获得相应公开题面，均标为 \texttt{external-textbook-statement-unavailable}。\end{solution}
\originalsection{第二部分}
\begin{exercise}\label{cw:ps9-II-34}原题 II.34：
(a) 生物学家研究楠塔基特岛的长尾鼬与草地田鼠。$x,y$ 是相对基线的数量，单位为百只：如 $x(2)=5$ 表示 $t=2$ 时多 500 只鼬，$y(2)=-5$ 表示少 500 只田鼠。模型为
\[x'=0.5x+y,\qquad y'=-2.25x+0.5y.\]
求特征值；取一根求复指数解，再取实、虚部产生两个实解。初始多 100 只鼬、田鼠恰为基线，求初值解；判断初始两种数量增减及田鼠何时再次等于基线。画约 $-2\le t\le2$ 的相图和两时间曲线，标明轨道上 $t=k\pi/3$、$k=-2,-1,0,1,2$ 五点，以说明三图对应。
(b) 使用上述 Mathlet，取消伴随矩阵选项、显示复平面特征值，将 $A=\begin{pmatrix}1&b\\0&1\end{pmatrix}$ 中 $b$ 从 $-4$ 调到 $4$。对每个 $b$ 核实特征值、描述全部特征向量、写全部正规模态 $e^{\lambda t}v$；若缺损则补一个独立解。各选一个可完全对角化与缺损的 $b$ 画相图。\end{exercise}
\begin{solution}
\textit{官方解译审与补证（AI 辅助）；其中另标编者补解或原文修正。}\par

(a) $\det(\lambda I-A)=(\lambda-1/2)^2+9/4$，根为 $\lambda=1/2\pm3i/2$。对正虚部根取 $v=(1,3i/2)^T$，则复解 $e^{(1/2+3i/2)t}v$ 的实、虚部为
\[\mathbf u_1=e^{t/2}\binom{\cos(3t/2)}{-(3/2)\sin(3t/2)},\qquad\mathbf u_2=e^{t/2}\binom{\sin(3t/2)}{(3/2)\cos(3t/2)}.\]
零点初值分别为 $(1,0)^T,(0,3/2)^T$，线性无关。给定初值 $\mathbf u(0)=(1,0)^T$ 恰选 $\mathbf u_1$。初始速度 $A(1,0)^T=(1/2,-9/4)^T$，因此鼬增加、田鼠减少，数量变化率分别为 $50$ 与 $-225$ 只每题面时间单位（原题未另指定时间单位）。$y(t)=0$ 在 $t=2n\pi/3$，故下一次为 $2\pi/3$。
相图为顺时针不稳定螺旋。在 $X=x,Y=2y/3$ 坐标中，$X^2+Y^2=e^t$ 且转角为 $-3t/2$，因此画图时形状、扩张及旋转方向均可直接检查。五个标记点为
\[\begin{array}{c|ccccc}
t&-2\pi/3&-\pi/3&0&\pi/3&2\pi/3\\\hline
(x,y)&(-e^{-\pi/3},0)&(0,\tfrac32e^{-\pi/6})&(1,0)&(0,-\tfrac32e^{\pi/6})&(-e^{\pi/3},0)
\end{array}\]
两端标记略超出 $[-2,2]$，所以图将时间范围延至 $[-2.1,2.1]$ 以保留原要求全部五点。
\begin{center}\begin{tikzpicture}
\begin{axis}[width=4.5cm,height=5cm,xmin=-3,xmax=2,ymin=-3,ymax=1.5,axis lines=middle,xlabel={$x$},ylabel={$y$},title={初值解的相轨道},ticks=none]
\addplot[main,thick,domain=-2.1:2.1,samples=160,->]({exp(x/2)*cos(deg(1.5*x))},{-1.5*exp(x/2)*sin(deg(1.5*x))});
\addplot[only marks,mark=*]coordinates{(-.35092,0)(0,.88858)(1,0)(0,-2.5321)(-2.84965,0)};
\node[below] at(axis cs:-.35092,0){$-2$};\node[above left]at(axis cs:0,.88858){$-1$};\node[above]at(axis cs:1,0){$0$};\node[right]at(axis cs:0,-2.5321){$1$};\node[above]at(axis cs:-2.84965,0){$2$};
\end{axis}
\begin{axis}[at={(4.6cm,0)},width=4.5cm,height=5cm,xmin=-2.1,xmax=2.1,ymin=-3,ymax=1.5,axis lines=middle,xlabel={$t$},title={$x(t)$},xtick={-2.0944,0,2.0944},xticklabels={$-2\pi/3$,$0$,$2\pi/3$},tick label style={font=\scriptsize}]
\addplot[main,thick,domain=-2.1:2.1,samples=150]{exp(x/2)*cos(deg(1.5*x))};
\end{axis}
\begin{axis}[at={(9.2cm,0)},width=4.5cm,height=5cm,xmin=-2.1,xmax=2.1,ymin=-3,ymax=1.5,axis lines=middle,xlabel={$t$},title={$y(t)$},xtick={-2.0944,0,2.0944},xticklabels={$-2\pi/3$,$0$,$2\pi/3$},tick label style={font=\scriptsize}]
\addplot[structurecolor,dashed,domain=-2.1:2.1,samples=150]{-1.5*exp(x/2)*sin(deg(1.5*x))};
\end{axis}\end{tikzpicture}\end{center}
轨道点旁的整数是 $k$，相应横坐标时刻由右图刻度给出；一个时刻的两曲线纵坐标正好组成左图一点。
(b) 这是按矩阵计算所得的核实，不是实际运行 Mathlet 的观察记录。特征多项式恒为 $(\lambda-1)^2$。若 $b=0$，矩阵为 $I$，所有非零向量均为特征向量，全部解都是 $e^tv$，形成径向向外的星形结点。若 $b\ne0$，$\ker(A-I)=\{(c,0)^T\}$，非零特征向量只在 $x$ 轴，正规模态为 $ce^t(1,0)^T$。取广义特征向量 $w=(0,1/b)^T$，满足 $(A-I)w=v=(1,0)^T$，补解为 $e^t(tv+w)$。等价的全部解为
\[\mathbf u(t)=e^t\binom{C_1+bC_2t}{C_2}.\]
其初值矩阵可逆，故已给齐全部解。$b=1$ 是缺损不稳定结点；在上半平面一般轨道向右剪切，角速度满足 $xy'-yx'=-by^2$，故 $b>0$ 为顺时针剪切，$b<0$ 为逆时针剪切。
\begin{center}\begin{tikzpicture}
\begin{axis}[width=5.8cm,height=4.8cm,xmin=-3,xmax=3,ymin=-3,ymax=3,axis lines=middle,xlabel={$x$},ylabel={$y$},title={$b=0$：不稳定星形结点},ticks=none]
\draw[->,main,thick](axis cs:0.30000,0.00000)--(axis cs:2.30000,0.00000);
\draw[->,main,thick](axis cs:0.21213,0.21213)--(axis cs:1.62635,1.62635);
\draw[->,main,thick](axis cs:0.00000,0.30000)--(axis cs:0.00000,2.30000);
\draw[->,main,thick](axis cs:-0.21213,0.21213)--(axis cs:-1.62635,1.62635);
\draw[->,main,thick](axis cs:-0.30000,0.00000)--(axis cs:-2.30000,0.00000);
\draw[->,main,thick](axis cs:-0.21213,-0.21213)--(axis cs:-1.62635,-1.62635);
\draw[->,main,thick](axis cs:-0.00000,-0.30000)--(axis cs:-0.00000,-2.30000);
\draw[->,main,thick](axis cs:0.21213,-0.21213)--(axis cs:1.62635,-1.62635);
\end{axis}
\begin{axis}[at={(6.1cm,0)},width=5.8cm,height=4.8cm,xmin=-3,xmax=3,ymin=-3,ymax=3,axis lines=middle,xlabel={$x$},ylabel={$y$},title={$b=1$：缺损不稳定结点},ticks=none]
\foreach \c/\d in {-1/1,1/1,-1/-1,1/-1}{\addplot[main,thick,domain=-3:.9,samples=90,->]({exp(x)*(\c+\d*x)},{exp(x)*\d});}
\draw[->,thick](axis cs:.2,0)--(axis cs:2.4,0);\draw[->,thick](axis cs:-.2,0)--(axis cs:-2.4,0);\node at(axis cs:1.3,-.45){特征线 $y=0$};
\end{axis}\end{tikzpicture}\end{center}
\end{solution}
\begin{exercise}\label{cw:ps9-II-35}原题 II.35：研究 $\mathbf u'=A\mathbf u$，$A=\begin{pmatrix}a&-3\\1&-1\end{pmatrix}$。在上述 Mathlet 中从 $a=-4$ 调到 $4$，关注迹--行列式平面的标记、复平面特征值及相图变化。
(a) 求迹、行列式及标记轨迹的方程。
(b) 求红边界的四个参数值，即 $\det A=0$、$\det A=(\operatorname{tr}A/2)^2$ 的两个值（一个超出 Mathlet 原滑条范围）及 $\operatorname{tr}A=0$。
(c) 在 $[-5,4]$ 上标四个值，分类五区间与四边界的九种相图；注明稳定性及有关旋转方向，分类可用螺旋、结点、鞍点、中心、星形、缺损结点、退化梳形等。
(d) 四个特殊值及五个区间各选一值画相图，注明所有特征线与时间方向。\end{exercise}
\begin{solution}
\textit{官方解译审与补证（AI 辅助）；其中另标编者补解或原文修正。}\par

(a) $T=\operatorname{tr}A=a-1$、$\Delta=\det A=3-a$，所以参数在 $(T,\Delta)$ 平面沿直线 $T+\Delta=2$ 运动。特征根为
\[\lambda_\pm=\frac{a-1\pm\sqrt{a^2+2a-11}}2.\]
(b) 四个分界为
\[a_-= -1-2\sqrt3\approx-4.4641,\quad1,\quad a_+=-1+2\sqrt3\approx2.4641,\quad3.\]
前后二次曲线交点来自 $4(3-a)=(a-1)^2$。$a_-$ 确在原滑条 $[-4,4]$ 外，须解析补足。
(c) 在复根区间，角速度的分子为
\[xy'-yx'=x^2-(a+1)xy+3y^2.\]
当 $a_-<a<a_+$ 时这是正定二次型，故均为逆时针。边界时退化为 $(x+\sqrt3y)^2$ 或 $(x-\sqrt3y)^2$，非特征线上的剪切亦为逆时针。分类如下；官方解页 2 有将 $a_+$ 的 $-1$ 写成 $+1$ 的笔误，此表按特征判别式核算。
\begin{center}\begin{tabular}{ll}\toprule
参数&类型\\\midrule
$a<a_-$&稳定结点\\
$a=a_-$&稳定缺损结点，逆时针剪切\\
$a_-<a<1$&逆时针稳定螺旋\\
$a=1$&逆时针中心\\
$1<a<a_+$&逆时针不稳定螺旋\\
$a=a_+$&不稳定缺损结点，逆时针剪切\\
$a_+<a<3$&不稳定结点\\
$a=3$&不稳定退化梳形\\
$a>3$&鞍点\\\bottomrule\end{tabular}\end{center}
不存在星形边界：在重根处矩阵不是 $\lambda I$，只有一条特征线。$a=3$ 时根为 $0,2$，$y=x$ 上全为平衡点，其他轨道沿 $(3,1)^T$ 的平行方向离开平衡线。
(d) 下列九图依次取 $a=-5,a_-,0,1,2,a_+,11/4,3,4$。虚线为特征线，实曲线为解析矩阵指数给出的代表轨道，曲线箭头或图内速度箭头标时间增加。任一实根的特征线为 $x=(\lambda+1)y$，这给出图中线的精确定位；中心图的轨道为闭椭圆。为可复算，作图使用
\[A=\frac T2I+B,\quad B^2=\frac{T^2-4\Delta}{4}I.\]
若右端系数为 $\rho^2>0$，则 $e^{At}=e^{Tt/2}[I\cosh\rho t+B\sinh(\rho t)/\rho]$；系数为 $-\beta^2<0$ 时改为 $I\cos\beta t+B\sin(\beta t)/\beta$；零时为 $e^{Tt/2}(I+tB)$。这由幂级数将奇、偶次幂分别求和得到。
\begin{center}
\begin{tikzpicture}
\begin{axis}[at={(0.0cm,0)},width=4.3cm,height=4cm,xmin=-2,xmax=2,ymin=-2,ymax=2,axis lines=middle,xlabel={$x$},ylabel={$y$},title={$a=-5$},ticks=none,clip=true]
\addplot[dashed] coordinates{(-2.2000,0.7333)(2.2000,-0.7333)};
\addplot[dashed] coordinates{(-2.2000,2.2000)(2.2000,-2.2000)};
\addplot[main,thin] coordinates{(2.2375,-0.4288)(1.9091,-0.3412)(1.6265,-0.2674)(1.3835,-0.2053)(1.1747,-0.1534)(0.9954,-0.1100)(0.8416,-0.0740)(0.7096,-0.0443)(0.5967,-0.0199)(0.5000,0.0000)(0.4174,0.0160)(0.3469,0.0289)(0.2868,0.0389)(0.2357,0.0467)(0.1922,0.0525)(0.1554,0.0568)(0.1243,0.0597)(0.0980,0.0615)(0.0759,0.0623)(0.0574,0.0625)(0.0419,0.0620)(0.0290,0.0611)(0.0183,0.0597)(0.0095,0.0581)(0.0024,0.0563)(-0.0035,0.0543)(-0.0081,0.0522)(-0.0118,0.0500)(-0.0147,0.0478)(-0.0168,0.0456)(-0.0184,0.0433)(-0.0196,0.0411)(-0.0203,0.0390)(-0.0207,0.0369)(-0.0208,0.0349)(-0.0207,0.0329)(-0.0205,0.0311)(-0.0201,0.0293)(-0.0196,0.0275)(-0.0190,0.0259)(-0.0184,0.0243)(-0.0177,0.0228)(-0.0169,0.0214)(-0.0162,0.0201)(-0.0155,0.0188)(-0.0147,0.0176)(-0.0140,0.0165)(-0.0133,0.0155)(-0.0126,0.0145)(-0.0119,0.0135)(-0.0112,0.0127)(-0.0106,0.0118)(-0.0100,0.0111)(-0.0094,0.0103)(-0.0088,0.0096)(-0.0083,0.0090)(-0.0078,0.0084)(-0.0073,0.0078)(-0.0069,0.0073)(-0.0064,0.0068)(-0.0060,0.0064)(-0.0056,0.0059)(-0.0053,0.0055)(-0.0049,0.0052)(-0.0046,0.0048)(-0.0043,0.0045)};
\addplot[main,thin] coordinates{(2.3972,0.3791)(1.9647,0.4421)(1.5975,0.4889)(1.2865,0.5221)(1.0237,0.5442)(0.8022,0.5570)(0.6160,0.5622)(0.4601,0.5612)(0.3301,0.5553)(0.2221,0.5455)(0.1329,0.5325)(0.0596,0.5171)(0.0000,0.5000)(-0.0481,0.4816)(-0.0866,0.4623)(-0.1168,0.4425)(-0.1401,0.4224)(-0.1576,0.4024)(-0.1703,0.3825)(-0.1790,0.3629)(-0.1844,0.3438)(-0.1870,0.3252)(-0.1874,0.3072)(-0.1860,0.2899)(-0.1832,0.2733)(-0.1792,0.2573)(-0.1744,0.2421)(-0.1689,0.2276)(-0.1629,0.2138)(-0.1566,0.2007)(-0.1500,0.1882)(-0.1434,0.1765)(-0.1367,0.1654)(-0.1300,0.1549)(-0.1234,0.1450)(-0.1170,0.1357)(-0.1107,0.1270)(-0.1047,0.1187)(-0.0988,0.1110)(-0.0932,0.1037)(-0.0878,0.0969)(-0.0826,0.0905)(-0.0777,0.0845)(-0.0730,0.0789)(-0.0685,0.0737)(-0.0643,0.0688)(-0.0603,0.0642)(-0.0565,0.0599)(-0.0529,0.0559)(-0.0496,0.0521)(-0.0464,0.0486)(-0.0434,0.0453)(-0.0406,0.0422)(-0.0380,0.0394)(-0.0355,0.0367)(-0.0332,0.0342)(-0.0310,0.0319)(-0.0289,0.0297)(-0.0270,0.0277)(-0.0252,0.0258)(-0.0235,0.0241)(-0.0220,0.0224)(-0.0205,0.0209)(-0.0191,0.0195)(-0.0178,0.0181)(-0.0166,0.0169)(-0.0155,0.0157)(-0.0145,0.0147)(-0.0135,0.0137)};
\addplot[main,thin] coordinates{(-2.2375,0.4288)(-1.9091,0.3412)(-1.6265,0.2674)(-1.3835,0.2053)(-1.1747,0.1534)(-0.9954,0.1100)(-0.8416,0.0740)(-0.7096,0.0443)(-0.5967,0.0199)(-0.5000,0.0000)(-0.4174,-0.0160)(-0.3469,-0.0289)(-0.2868,-0.0389)(-0.2357,-0.0467)(-0.1922,-0.0525)(-0.1554,-0.0568)(-0.1243,-0.0597)(-0.0980,-0.0615)(-0.0759,-0.0623)(-0.0574,-0.0625)(-0.0419,-0.0620)(-0.0290,-0.0611)(-0.0183,-0.0597)(-0.0095,-0.0581)(-0.0024,-0.0563)(0.0035,-0.0543)(0.0081,-0.0522)(0.0118,-0.0500)(0.0147,-0.0478)(0.0168,-0.0456)(0.0184,-0.0433)(0.0196,-0.0411)(0.0203,-0.0390)(0.0207,-0.0369)(0.0208,-0.0349)(0.0207,-0.0329)(0.0205,-0.0311)(0.0201,-0.0293)(0.0196,-0.0275)(0.0190,-0.0259)(0.0184,-0.0243)(0.0177,-0.0228)(0.0169,-0.0214)(0.0162,-0.0201)(0.0155,-0.0188)(0.0147,-0.0176)(0.0140,-0.0165)(0.0133,-0.0155)(0.0126,-0.0145)(0.0119,-0.0135)(0.0112,-0.0127)(0.0106,-0.0118)(0.0100,-0.0111)(0.0094,-0.0103)(0.0088,-0.0096)(0.0083,-0.0090)(0.0078,-0.0084)(0.0073,-0.0078)(0.0069,-0.0073)(0.0064,-0.0068)(0.0060,-0.0064)(0.0056,-0.0059)(0.0053,-0.0055)(0.0049,-0.0052)(0.0046,-0.0048)(0.0043,-0.0045)};
\addplot[main,thin] coordinates{(-2.3972,-0.3791)(-1.9647,-0.4421)(-1.5975,-0.4889)(-1.2865,-0.5221)(-1.0237,-0.5442)(-0.8022,-0.5570)(-0.6160,-0.5622)(-0.4601,-0.5612)(-0.3301,-0.5553)(-0.2221,-0.5455)(-0.1329,-0.5325)(-0.0596,-0.5171)(0.0000,-0.5000)(0.0481,-0.4816)(0.0866,-0.4623)(0.1168,-0.4425)(0.1401,-0.4224)(0.1576,-0.4024)(0.1703,-0.3825)(0.1790,-0.3629)(0.1844,-0.3438)(0.1870,-0.3252)(0.1874,-0.3072)(0.1860,-0.2899)(0.1832,-0.2733)(0.1792,-0.2573)(0.1744,-0.2421)(0.1689,-0.2276)(0.1629,-0.2138)(0.1566,-0.2007)(0.1500,-0.1882)(0.1434,-0.1765)(0.1367,-0.1654)(0.1300,-0.1549)(0.1234,-0.1450)(0.1170,-0.1357)(0.1107,-0.1270)(0.1047,-0.1187)(0.0988,-0.1110)(0.0932,-0.1037)(0.0878,-0.0969)(0.0826,-0.0905)(0.0777,-0.0845)(0.0730,-0.0789)(0.0685,-0.0737)(0.0643,-0.0688)(0.0603,-0.0642)(0.0565,-0.0599)(0.0529,-0.0559)(0.0496,-0.0521)(0.0464,-0.0486)(0.0434,-0.0453)(0.0406,-0.0422)(0.0380,-0.0394)(0.0355,-0.0367)(0.0332,-0.0342)(0.0310,-0.0319)(0.0289,-0.0297)(0.0270,-0.0277)(0.0252,-0.0258)(0.0235,-0.0241)(0.0220,-0.0224)(0.0205,-0.0209)(0.0191,-0.0195)(0.0178,-0.0181)(0.0166,-0.0169)(0.0155,-0.0157)(0.0145,-0.0147)(0.0135,-0.0137)};
\draw[->,thick](axis cs:1.0,0.0)--(axis cs:0.6862,0.0628);
\draw[->,thick](axis cs:0.0,1.0)--(axis cs:-0.3036,0.8988);
\draw[->,thick](axis cs:-1.0,0.0)--(axis cs:-0.6862,-0.0628);
\draw[->,thick](axis cs:0.0,-1.0)--(axis cs:0.3036,-0.8988);
\end{axis}
\begin{axis}[at={(4.35cm,0)},width=4.3cm,height=4cm,xmin=-2,xmax=2,ymin=-2,ymax=2,axis lines=middle,xlabel={$x$},ylabel={$y$},title={$a=a_-$},ticks=none,clip=true]
\addplot[dashed] coordinates{(-2.2000,1.2702)(2.2000,-1.2702)};
\addplot[dashed] coordinates{(-2.2000,1.2702)(2.2000,-1.2702)};
\addplot[main,thin] coordinates{(2.4577,-0.5746)(2.1471,-0.4738)(1.8731,-0.3868)(1.6315,-0.3118)(1.4186,-0.2475)(1.2311,-0.1924)(1.0663,-0.1454)(0.9215,-0.1055)(0.7944,-0.0718)(0.6829,-0.0434)(0.5853,-0.0197)(0.5000,0.0000)(0.4255,0.0162)(0.3605,0.0294)(0.3039,0.0400)(0.2547,0.0483)(0.2120,0.0548)(0.1751,0.0597)(0.1432,0.0631)(0.1157,0.0654)(0.0921,0.0668)(0.0719,0.0673)(0.0546,0.0672)(0.0400,0.0664)(0.0275,0.0653)(0.0171,0.0638)(0.0083,0.0620)(0.0011,0.0600)(-0.0049,0.0578)(-0.0098,0.0555)(-0.0137,0.0531)(-0.0168,0.0507)(-0.0193,0.0483)(-0.0211,0.0459)(-0.0224,0.0435)(-0.0233,0.0412)(-0.0238,0.0389)(-0.0241,0.0367)(-0.0240,0.0346)(-0.0238,0.0325)(-0.0234,0.0306)(-0.0229,0.0287)(-0.0223,0.0269)(-0.0216,0.0252)(-0.0208,0.0235)(-0.0200,0.0220)(-0.0192,0.0205)(-0.0183,0.0192)(-0.0174,0.0179)(-0.0166,0.0166)(-0.0157,0.0155)(-0.0149,0.0144)(-0.0141,0.0134)(-0.0133,0.0125)(-0.0125,0.0116)(-0.0118,0.0107)(-0.0111,0.0100)(-0.0104,0.0092)(-0.0097,0.0086)(-0.0091,0.0079)(-0.0085,0.0073)(-0.0080,0.0068)(-0.0074,0.0063)(-0.0069,0.0058)(-0.0065,0.0054)(-0.0060,0.0050)(-0.0056,0.0046)(-0.0052,0.0042)};
\addplot[main,thin] coordinates{(2.4760,0.3481)(2.0731,0.4155)(1.7237,0.4673)(1.4213,0.5060)(1.1603,0.5334)(0.9355,0.5513)(0.7424,0.5613)(0.5772,0.5646)(0.4363,0.5625)(0.3166,0.5559)(0.2154,0.5457)(0.1302,0.5326)(0.0591,0.5171)(0.0000,0.5000)(-0.0486,0.4816)(-0.0881,0.4622)(-0.1199,0.4424)(-0.1450,0.4222)(-0.1644,0.4019)(-0.1790,0.3818)(-0.1894,0.3619)(-0.1963,0.3424)(-0.2004,0.3234)(-0.2019,0.3050)(-0.2015,0.2873)(-0.1993,0.2701)(-0.1959,0.2537)(-0.1913,0.2380)(-0.1859,0.2231)(-0.1799,0.2088)(-0.1734,0.1953)(-0.1665,0.1825)(-0.1594,0.1704)(-0.1522,0.1589)(-0.1450,0.1481)(-0.1378,0.1380)(-0.1306,0.1284)(-0.1236,0.1195)(-0.1168,0.1111)(-0.1102,0.1032)(-0.1038,0.0958)(-0.0976,0.0889)(-0.0917,0.0825)(-0.0861,0.0765)(-0.0807,0.0709)(-0.0755,0.0656)(-0.0706,0.0608)(-0.0660,0.0562)(-0.0616,0.0520)(-0.0575,0.0481)(-0.0536,0.0445)(-0.0499,0.0411)(-0.0465,0.0380)(-0.0432,0.0351)(-0.0402,0.0324)(-0.0374,0.0299)(-0.0347,0.0276)(-0.0322,0.0254)(-0.0299,0.0234)(-0.0277,0.0216)(-0.0257,0.0199)(-0.0238,0.0184)(-0.0220,0.0169)(-0.0204,0.0156)(-0.0189,0.0143)(-0.0174,0.0132)(-0.0161,0.0121)(-0.0149,0.0112)(-0.0138,0.0103)(-0.0127,0.0095)};
\addplot[main,thin] coordinates{(-2.4577,0.5746)(-2.1471,0.4738)(-1.8731,0.3868)(-1.6315,0.3118)(-1.4186,0.2475)(-1.2311,0.1924)(-1.0663,0.1454)(-0.9215,0.1055)(-0.7944,0.0718)(-0.6829,0.0434)(-0.5853,0.0197)(-0.5000,0.0000)(-0.4255,-0.0162)(-0.3605,-0.0294)(-0.3039,-0.0400)(-0.2547,-0.0483)(-0.2120,-0.0548)(-0.1751,-0.0597)(-0.1432,-0.0631)(-0.1157,-0.0654)(-0.0921,-0.0668)(-0.0719,-0.0673)(-0.0546,-0.0672)(-0.0400,-0.0664)(-0.0275,-0.0653)(-0.0171,-0.0638)(-0.0083,-0.0620)(-0.0011,-0.0600)(0.0049,-0.0578)(0.0098,-0.0555)(0.0137,-0.0531)(0.0168,-0.0507)(0.0193,-0.0483)(0.0211,-0.0459)(0.0224,-0.0435)(0.0233,-0.0412)(0.0238,-0.0389)(0.0241,-0.0367)(0.0240,-0.0346)(0.0238,-0.0325)(0.0234,-0.0306)(0.0229,-0.0287)(0.0223,-0.0269)(0.0216,-0.0252)(0.0208,-0.0235)(0.0200,-0.0220)(0.0192,-0.0205)(0.0183,-0.0192)(0.0174,-0.0179)(0.0166,-0.0166)(0.0157,-0.0155)(0.0149,-0.0144)(0.0141,-0.0134)(0.0133,-0.0125)(0.0125,-0.0116)(0.0118,-0.0107)(0.0111,-0.0100)(0.0104,-0.0092)(0.0097,-0.0086)(0.0091,-0.0079)(0.0085,-0.0073)(0.0080,-0.0068)(0.0074,-0.0063)(0.0069,-0.0058)(0.0065,-0.0054)(0.0060,-0.0050)(0.0056,-0.0046)(0.0052,-0.0042)};
\addplot[main,thin] coordinates{(-2.4760,-0.3481)(-2.0731,-0.4155)(-1.7237,-0.4673)(-1.4213,-0.5060)(-1.1603,-0.5334)(-0.9355,-0.5513)(-0.7424,-0.5613)(-0.5772,-0.5646)(-0.4363,-0.5625)(-0.3166,-0.5559)(-0.2154,-0.5457)(-0.1302,-0.5326)(-0.0591,-0.5171)(0.0000,-0.5000)(0.0486,-0.4816)(0.0881,-0.4622)(0.1199,-0.4424)(0.1450,-0.4222)(0.1644,-0.4019)(0.1790,-0.3818)(0.1894,-0.3619)(0.1963,-0.3424)(0.2004,-0.3234)(0.2019,-0.3050)(0.2015,-0.2873)(0.1993,-0.2701)(0.1959,-0.2537)(0.1913,-0.2380)(0.1859,-0.2231)(0.1799,-0.2088)(0.1734,-0.1953)(0.1665,-0.1825)(0.1594,-0.1704)(0.1522,-0.1589)(0.1450,-0.1481)(0.1378,-0.1380)(0.1306,-0.1284)(0.1236,-0.1195)(0.1168,-0.1111)(0.1102,-0.1032)(0.1038,-0.0958)(0.0976,-0.0889)(0.0917,-0.0825)(0.0861,-0.0765)(0.0807,-0.0709)(0.0755,-0.0656)(0.0706,-0.0608)(0.0660,-0.0562)(0.0616,-0.0520)(0.0575,-0.0481)(0.0536,-0.0445)(0.0499,-0.0411)(0.0465,-0.0380)(0.0432,-0.0351)(0.0402,-0.0324)(0.0374,-0.0299)(0.0347,-0.0276)(0.0322,-0.0254)(0.0299,-0.0234)(0.0277,-0.0216)(0.0257,-0.0199)(0.0238,-0.0184)(0.0220,-0.0169)(0.0204,-0.0156)(0.0189,-0.0143)(0.0174,-0.0132)(0.0161,-0.0121)(0.0149,-0.0112)(0.0138,-0.0103)(0.0127,-0.0095)};
\draw[->,thick](axis cs:1.0,0.0)--(axis cs:0.6877,0.0699);
\draw[->,thick](axis cs:0.0,1.0)--(axis cs:-0.3036,0.8988);
\draw[->,thick](axis cs:-1.0,0.0)--(axis cs:-0.6877,-0.0699);
\draw[->,thick](axis cs:0.0,-1.0)--(axis cs:0.3036,-0.8988);
\end{axis}
\begin{axis}[at={(8.7cm,0)},width=4.3cm,height=4cm,xmin=-2,xmax=2,ymin=-2,ymax=2,axis lines=middle,xlabel={$x$},ylabel={$y$},title={$a=0$},ticks=none,clip=true]
\addplot[main,thin] coordinates{(-0.3197,-0.4952)(-0.2670,-0.4881)(-0.2151,-0.4794)(-0.1643,-0.4693)(-0.1146,-0.4577)(-0.0663,-0.4448)(-0.0193,-0.4307)(0.0260,-0.4154)(0.0696,-0.3992)(0.1115,-0.3820)(0.1515,-0.3640)(0.1895,-0.3452)(0.2254,-0.3258)(0.2593,-0.3059)(0.2909,-0.2855)(0.3204,-0.2647)(0.3477,-0.2437)(0.3726,-0.2225)(0.3953,-0.2012)(0.4158,-0.1799)(0.4339,-0.1587)(0.4498,-0.1376)(0.4634,-0.1168)(0.4748,-0.0962)(0.4840,-0.0760)(0.4911,-0.0562)(0.4961,-0.0369)(0.4990,-0.0182)(0.5000,0.0000)(0.4991,0.0175)(0.4963,0.0344)(0.4917,0.0505)(0.4855,0.0659)(0.4776,0.0805)(0.4683,0.0942)(0.4575,0.1072)(0.4453,0.1193)(0.4319,0.1305)(0.4174,0.1408)(0.4018,0.1502)(0.3852,0.1588)(0.3678,0.1664)(0.3496,0.1731)(0.3307,0.1790)(0.3113,0.1840)(0.2913,0.1881)(0.2710,0.1914)(0.2504,0.1938)(0.2295,0.1954)(0.2085,0.1962)(0.1875,0.1963)(0.1665,0.1956)(0.1456,0.1942)(0.1249,0.1922)(0.1044,0.1894)(0.0843,0.1861)(0.0646,0.1822)(0.0453,0.1777)(0.0265,0.1727)(0.0083,0.1673)(-0.0093,0.1614)(-0.0263,0.1551)(-0.0425,0.1484)(-0.0581,0.1415)(-0.0728,0.1342)(-0.0868,0.1267)(-0.1000,0.1190)(-0.1123,0.1111)(-0.1238,0.1030)(-0.1344,0.0949)(-0.1441,0.0867)(-0.1529,0.0784)(-0.1609,0.0701)(-0.1680,0.0619)(-0.1742,0.0537)(-0.1795,0.0456)(-0.1840,0.0377)(-0.1876,0.0298)(-0.1904,0.0221)(-0.1923,0.0147)(-0.1935,0.0074)(-0.1939,0.0003)(-0.1936,-0.0065)(-0.1925,-0.0130)(-0.1908,-0.0193)};
\addplot[main,thin] coordinates{(1.4856,0.1756)(1.4643,0.2212)(1.4383,0.2643)(1.4078,0.3050)(1.3731,0.3431)(1.3344,0.3785)(1.2920,0.4113)(1.2463,0.4414)(1.1975,0.4688)(1.1460,0.4935)(1.0919,0.5154)(1.0356,0.5347)(0.9774,0.5512)(0.9176,0.5651)(0.8564,0.5764)(0.7942,0.5851)(0.7311,0.5914)(0.6675,0.5951)(0.6037,0.5966)(0.5398,0.5957)(0.4761,0.5926)(0.4129,0.5874)(0.3503,0.5802)(0.2886,0.5710)(0.2280,0.5600)(0.1687,0.5473)(0.1108,0.5330)(0.0545,0.5172)(0.0000,0.5000)(-0.0526,0.4815)(-0.1031,0.4619)(-0.1515,0.4412)(-0.1977,0.4196)(-0.2414,0.3972)(-0.2827,0.3740)(-0.3215,0.3503)(-0.3578,0.3261)(-0.3914,0.3015)(-0.4224,0.2766)(-0.4507,0.2516)(-0.4763,0.2265)(-0.4992,0.2014)(-0.5194,0.1765)(-0.5370,0.1517)(-0.5520,0.1273)(-0.5643,0.1032)(-0.5741,0.0796)(-0.5814,0.0566)(-0.5862,0.0341)(-0.5887,0.0123)(-0.5889,-0.0088)(-0.5869,-0.0291)(-0.5827,-0.0486)(-0.5765,-0.0673)(-0.5683,-0.0850)(-0.5583,-0.1018)(-0.5465,-0.1176)(-0.5331,-0.1324)(-0.5182,-0.1462)(-0.5018,-0.1590)(-0.4842,-0.1707)(-0.4653,-0.1814)(-0.4453,-0.1910)(-0.4244,-0.1995)(-0.4026,-0.2070)(-0.3801,-0.2135)(-0.3569,-0.2189)(-0.3332,-0.2234)(-0.3091,-0.2268)(-0.2846,-0.2293)(-0.2600,-0.2308)(-0.2352,-0.2313)(-0.2104,-0.2310)(-0.1857,-0.2299)(-0.1612,-0.2279)(-0.1369,-0.2251)(-0.1130,-0.2216)(-0.0895,-0.2174)(-0.0664,-0.2125)(-0.0440,-0.2070)(-0.0221,-0.2009)(-0.0009,-0.1942)(0.0195,-0.1871)(0.0391,-0.1795)(0.0579,-0.1715)};
\addplot[main,thin] coordinates{(0.3197,0.4952)(0.2670,0.4881)(0.2151,0.4794)(0.1643,0.4693)(0.1146,0.4577)(0.0663,0.4448)(0.0193,0.4307)(-0.0260,0.4154)(-0.0696,0.3992)(-0.1115,0.3820)(-0.1515,0.3640)(-0.1895,0.3452)(-0.2254,0.3258)(-0.2593,0.3059)(-0.2909,0.2855)(-0.3204,0.2647)(-0.3477,0.2437)(-0.3726,0.2225)(-0.3953,0.2012)(-0.4158,0.1799)(-0.4339,0.1587)(-0.4498,0.1376)(-0.4634,0.1168)(-0.4748,0.0962)(-0.4840,0.0760)(-0.4911,0.0562)(-0.4961,0.0369)(-0.4990,0.0182)(-0.5000,0.0000)(-0.4991,-0.0175)(-0.4963,-0.0344)(-0.4917,-0.0505)(-0.4855,-0.0659)(-0.4776,-0.0805)(-0.4683,-0.0942)(-0.4575,-0.1072)(-0.4453,-0.1193)(-0.4319,-0.1305)(-0.4174,-0.1408)(-0.4018,-0.1502)(-0.3852,-0.1588)(-0.3678,-0.1664)(-0.3496,-0.1731)(-0.3307,-0.1790)(-0.3113,-0.1840)(-0.2913,-0.1881)(-0.2710,-0.1914)(-0.2504,-0.1938)(-0.2295,-0.1954)(-0.2085,-0.1962)(-0.1875,-0.1963)(-0.1665,-0.1956)(-0.1456,-0.1942)(-0.1249,-0.1922)(-0.1044,-0.1894)(-0.0843,-0.1861)(-0.0646,-0.1822)(-0.0453,-0.1777)(-0.0265,-0.1727)(-0.0083,-0.1673)(0.0093,-0.1614)(0.0263,-0.1551)(0.0425,-0.1484)(0.0581,-0.1415)(0.0728,-0.1342)(0.0868,-0.1267)(0.1000,-0.1190)(0.1123,-0.1111)(0.1238,-0.1030)(0.1344,-0.0949)(0.1441,-0.0867)(0.1529,-0.0784)(0.1609,-0.0701)(0.1680,-0.0619)(0.1742,-0.0537)(0.1795,-0.0456)(0.1840,-0.0377)(0.1876,-0.0298)(0.1904,-0.0221)(0.1923,-0.0147)(0.1935,-0.0074)(0.1939,-0.0003)(0.1936,0.0065)(0.1925,0.0130)(0.1908,0.0193)};
\addplot[main,thin] coordinates{(-1.4856,-0.1756)(-1.4643,-0.2212)(-1.4383,-0.2643)(-1.4078,-0.3050)(-1.3731,-0.3431)(-1.3344,-0.3785)(-1.2920,-0.4113)(-1.2463,-0.4414)(-1.1975,-0.4688)(-1.1460,-0.4935)(-1.0919,-0.5154)(-1.0356,-0.5347)(-0.9774,-0.5512)(-0.9176,-0.5651)(-0.8564,-0.5764)(-0.7942,-0.5851)(-0.7311,-0.5914)(-0.6675,-0.5951)(-0.6037,-0.5966)(-0.5398,-0.5957)(-0.4761,-0.5926)(-0.4129,-0.5874)(-0.3503,-0.5802)(-0.2886,-0.5710)(-0.2280,-0.5600)(-0.1687,-0.5473)(-0.1108,-0.5330)(-0.0545,-0.5172)(0.0000,-0.5000)(0.0526,-0.4815)(0.1031,-0.4619)(0.1515,-0.4412)(0.1977,-0.4196)(0.2414,-0.3972)(0.2827,-0.3740)(0.3215,-0.3503)(0.3578,-0.3261)(0.3914,-0.3015)(0.4224,-0.2766)(0.4507,-0.2516)(0.4763,-0.2265)(0.4992,-0.2014)(0.5194,-0.1765)(0.5370,-0.1517)(0.5520,-0.1273)(0.5643,-0.1032)(0.5741,-0.0796)(0.5814,-0.0566)(0.5862,-0.0341)(0.5887,-0.0123)(0.5889,0.0088)(0.5869,0.0291)(0.5827,0.0486)(0.5765,0.0673)(0.5683,0.0850)(0.5583,0.1018)(0.5465,0.1176)(0.5331,0.1324)(0.5182,0.1462)(0.5018,0.1590)(0.4842,0.1707)(0.4653,0.1814)(0.4453,0.1910)(0.4244,0.1995)(0.4026,0.2070)(0.3801,0.2135)(0.3569,0.2189)(0.3332,0.2234)(0.3091,0.2268)(0.2846,0.2293)(0.2600,0.2308)(0.2352,0.2313)(0.2104,0.2310)(0.1857,0.2299)(0.1612,0.2279)(0.1369,0.2251)(0.1130,0.2216)(0.0895,0.2174)(0.0664,0.2125)(0.0440,0.2070)(0.0221,0.2009)(0.0009,0.1942)(-0.0195,0.1871)(-0.0391,0.1795)(-0.0579,0.1715)};
\draw[->,thick](axis cs:1.0,0.0)--(axis cs:1.0000,0.3200);
\draw[->,thick](axis cs:0.0,1.0)--(axis cs:-0.3036,0.8988);
\draw[->,thick](axis cs:-1.0,0.0)--(axis cs:-1.0000,-0.3200);
\draw[->,thick](axis cs:0.0,-1.0)--(axis cs:0.3036,-0.8988);
\end{axis}
\end{tikzpicture}
\par
\begin{tikzpicture}
\begin{axis}[at={(0.0cm,0)},width=4.3cm,height=4cm,xmin=-2,xmax=2,ymin=-2,ymax=2,axis lines=middle,xlabel={$x$},ylabel={$y$},title={$a=1$},ticks=none,clip=true]
\addplot[main,thin] coordinates{(0.5000,0.0000)(0.5214,0.0224)(0.5407,0.0448)(0.5579,0.0669)(0.5728,0.0888)(0.5854,0.1103)(0.5956,0.1314)(0.6034,0.1520)(0.6088,0.1719)(0.6118,0.1911)(0.6123,0.2096)(0.6103,0.2273)(0.6058,0.2440)(0.5990,0.2597)(0.5897,0.2744)(0.5780,0.2880)(0.5640,0.3004)(0.5478,0.3116)(0.5293,0.3216)(0.5087,0.3303)(0.4861,0.3376)(0.4615,0.3436)(0.4350,0.3482)(0.4068,0.3514)(0.3769,0.3532)(0.3456,0.3535)(0.3128,0.3524)(0.2788,0.3500)(0.2437,0.3461)(0.2075,0.3408)(0.1706,0.3341)(0.1329,0.3261)(0.0947,0.3168)(0.0562,0.3062)(0.0174,0.2944)(-0.0215,0.2813)(-0.0602,0.2672)(-0.0988,0.2520)(-0.1369,0.2357)(-0.1745,0.2186)(-0.2114,0.2005)(-0.2474,0.1816)(-0.2824,0.1620)(-0.3163,0.1418)(-0.3489,0.1209)(-0.3801,0.0996)(-0.4098,0.0779)(-0.4379,0.0559)(-0.4641,0.0336)(-0.4885,0.0112)(-0.5110,-0.0112)(-0.5313,-0.0336)(-0.5496,-0.0559)(-0.5656,-0.0779)(-0.5794,-0.0996)(-0.5908,-0.1209)(-0.5998,-0.1418)(-0.6064,-0.1620)(-0.6106,-0.1816)(-0.6123,-0.2005)(-0.6116,-0.2186)(-0.6084,-0.2357)(-0.6027,-0.2520)(-0.5946,-0.2672)(-0.5841,-0.2813)(-0.5713,-0.2944)(-0.5562,-0.3062)(-0.5388,-0.3168)(-0.5193,-0.3261)(-0.4976,-0.3341)(-0.4740,-0.3408)(-0.4485,-0.3461)(-0.4211,-0.3500)(-0.3921,-0.3524)(-0.3614,-0.3535)(-0.3294,-0.3532)(-0.2960,-0.3514)(-0.2614,-0.3482)(-0.2257,-0.3436)(-0.1891,-0.3376)(-0.1518,-0.3303)(-0.1139,-0.3216)(-0.0755,-0.3116)(-0.0368,-0.3004)(0.0020,-0.2880)(0.0409,-0.2744)(0.0795,-0.2597)(0.1179,-0.2440)(0.1558,-0.2273)(0.1930,-0.2096)(0.2295,-0.1911)(0.2650,-0.1719)(0.2995,-0.1520)(0.3328,-0.1314)(0.3647,-0.1103)(0.3952,-0.0888)(0.4241,-0.0669)(0.4512,-0.0448)(0.4766,-0.0224)(0.5000,-0.0000)};
\addplot[main,thin] coordinates{(0.0000,0.5000)(-0.0673,0.4766)(-0.1343,0.4512)(-0.2007,0.4241)(-0.2664,0.3952)(-0.3310,0.3647)(-0.3942,0.3328)(-0.4559,0.2995)(-0.5157,0.2650)(-0.5734,0.2295)(-0.6289,0.1930)(-0.6818,0.1558)(-0.7319,0.1179)(-0.7792,0.0795)(-0.8232,0.0409)(-0.8640,0.0020)(-0.9013,-0.0368)(-0.9349,-0.0755)(-0.9648,-0.1139)(-0.9908,-0.1518)(-1.0128,-0.1891)(-1.0308,-0.2257)(-1.0445,-0.2614)(-1.0541,-0.2960)(-1.0595,-0.3294)(-1.0605,-0.3614)(-1.0573,-0.3921)(-1.0499,-0.4211)(-1.0382,-0.4485)(-1.0223,-0.4740)(-1.0023,-0.4976)(-0.9783,-0.5193)(-0.9503,-0.5388)(-0.9186,-0.5562)(-0.8831,-0.5713)(-0.8440,-0.5841)(-0.8016,-0.5946)(-0.7559,-0.6027)(-0.7072,-0.6084)(-0.6557,-0.6116)(-0.6015,-0.6123)(-0.5448,-0.6106)(-0.4860,-0.6064)(-0.4253,-0.5998)(-0.3628,-0.5908)(-0.2988,-0.5794)(-0.2337,-0.5656)(-0.1676,-0.5496)(-0.1008,-0.5313)(-0.0337,-0.5110)(0.0337,-0.4885)(0.1008,-0.4641)(0.1676,-0.4379)(0.2337,-0.4098)(0.2988,-0.3801)(0.3628,-0.3489)(0.4253,-0.3163)(0.4860,-0.2824)(0.5448,-0.2474)(0.6015,-0.2114)(0.6557,-0.1745)(0.7072,-0.1369)(0.7559,-0.0988)(0.8016,-0.0602)(0.8440,-0.0215)(0.8831,0.0174)(0.9186,0.0562)(0.9503,0.0947)(0.9783,0.1329)(1.0023,0.1706)(1.0223,0.2075)(1.0382,0.2437)(1.0499,0.2788)(1.0573,0.3128)(1.0605,0.3456)(1.0595,0.3769)(1.0541,0.4068)(1.0445,0.4350)(1.0308,0.4615)(1.0128,0.4861)(0.9908,0.5087)(0.9648,0.5293)(0.9349,0.5478)(0.9013,0.5640)(0.8640,0.5780)(0.8232,0.5897)(0.7792,0.5990)(0.7319,0.6058)(0.6818,0.6103)(0.6289,0.6123)(0.5734,0.6118)(0.5157,0.6088)(0.4559,0.6034)(0.3942,0.5956)(0.3310,0.5854)(0.2664,0.5728)(0.2007,0.5579)(0.1343,0.5407)(0.0673,0.5214)(-0.0000,0.5000)};
\addplot[main,thin] coordinates{(-0.5000,0.0000)(-0.5214,-0.0224)(-0.5407,-0.0448)(-0.5579,-0.0669)(-0.5728,-0.0888)(-0.5854,-0.1103)(-0.5956,-0.1314)(-0.6034,-0.1520)(-0.6088,-0.1719)(-0.6118,-0.1911)(-0.6123,-0.2096)(-0.6103,-0.2273)(-0.6058,-0.2440)(-0.5990,-0.2597)(-0.5897,-0.2744)(-0.5780,-0.2880)(-0.5640,-0.3004)(-0.5478,-0.3116)(-0.5293,-0.3216)(-0.5087,-0.3303)(-0.4861,-0.3376)(-0.4615,-0.3436)(-0.4350,-0.3482)(-0.4068,-0.3514)(-0.3769,-0.3532)(-0.3456,-0.3535)(-0.3128,-0.3524)(-0.2788,-0.3500)(-0.2437,-0.3461)(-0.2075,-0.3408)(-0.1706,-0.3341)(-0.1329,-0.3261)(-0.0947,-0.3168)(-0.0562,-0.3062)(-0.0174,-0.2944)(0.0215,-0.2813)(0.0602,-0.2672)(0.0988,-0.2520)(0.1369,-0.2357)(0.1745,-0.2186)(0.2114,-0.2005)(0.2474,-0.1816)(0.2824,-0.1620)(0.3163,-0.1418)(0.3489,-0.1209)(0.3801,-0.0996)(0.4098,-0.0779)(0.4379,-0.0559)(0.4641,-0.0336)(0.4885,-0.0112)(0.5110,0.0112)(0.5313,0.0336)(0.5496,0.0559)(0.5656,0.0779)(0.5794,0.0996)(0.5908,0.1209)(0.5998,0.1418)(0.6064,0.1620)(0.6106,0.1816)(0.6123,0.2005)(0.6116,0.2186)(0.6084,0.2357)(0.6027,0.2520)(0.5946,0.2672)(0.5841,0.2813)(0.5713,0.2944)(0.5562,0.3062)(0.5388,0.3168)(0.5193,0.3261)(0.4976,0.3341)(0.4740,0.3408)(0.4485,0.3461)(0.4211,0.3500)(0.3921,0.3524)(0.3614,0.3535)(0.3294,0.3532)(0.2960,0.3514)(0.2614,0.3482)(0.2257,0.3436)(0.1891,0.3376)(0.1518,0.3303)(0.1139,0.3216)(0.0755,0.3116)(0.0368,0.3004)(-0.0020,0.2880)(-0.0409,0.2744)(-0.0795,0.2597)(-0.1179,0.2440)(-0.1558,0.2273)(-0.1930,0.2096)(-0.2295,0.1911)(-0.2650,0.1719)(-0.2995,0.1520)(-0.3328,0.1314)(-0.3647,0.1103)(-0.3952,0.0888)(-0.4241,0.0669)(-0.4512,0.0448)(-0.4766,0.0224)(-0.5000,0.0000)};
\addplot[main,thin] coordinates{(0.0000,-0.5000)(0.0673,-0.4766)(0.1343,-0.4512)(0.2007,-0.4241)(0.2664,-0.3952)(0.3310,-0.3647)(0.3942,-0.3328)(0.4559,-0.2995)(0.5157,-0.2650)(0.5734,-0.2295)(0.6289,-0.1930)(0.6818,-0.1558)(0.7319,-0.1179)(0.7792,-0.0795)(0.8232,-0.0409)(0.8640,-0.0020)(0.9013,0.0368)(0.9349,0.0755)(0.9648,0.1139)(0.9908,0.1518)(1.0128,0.1891)(1.0308,0.2257)(1.0445,0.2614)(1.0541,0.2960)(1.0595,0.3294)(1.0605,0.3614)(1.0573,0.3921)(1.0499,0.4211)(1.0382,0.4485)(1.0223,0.4740)(1.0023,0.4976)(0.9783,0.5193)(0.9503,0.5388)(0.9186,0.5562)(0.8831,0.5713)(0.8440,0.5841)(0.8016,0.5946)(0.7559,0.6027)(0.7072,0.6084)(0.6557,0.6116)(0.6015,0.6123)(0.5448,0.6106)(0.4860,0.6064)(0.4253,0.5998)(0.3628,0.5908)(0.2988,0.5794)(0.2337,0.5656)(0.1676,0.5496)(0.1008,0.5313)(0.0337,0.5110)(-0.0337,0.4885)(-0.1008,0.4641)(-0.1676,0.4379)(-0.2337,0.4098)(-0.2988,0.3801)(-0.3628,0.3489)(-0.4253,0.3163)(-0.4860,0.2824)(-0.5448,0.2474)(-0.6015,0.2114)(-0.6557,0.1745)(-0.7072,0.1369)(-0.7559,0.0988)(-0.8016,0.0602)(-0.8440,0.0215)(-0.8831,-0.0174)(-0.9186,-0.0562)(-0.9503,-0.0947)(-0.9783,-0.1329)(-1.0023,-0.1706)(-1.0223,-0.2075)(-1.0382,-0.2437)(-1.0499,-0.2788)(-1.0573,-0.3128)(-1.0605,-0.3456)(-1.0595,-0.3769)(-1.0541,-0.4068)(-1.0445,-0.4350)(-1.0308,-0.4615)(-1.0128,-0.4861)(-0.9908,-0.5087)(-0.9648,-0.5293)(-0.9349,-0.5478)(-0.9013,-0.5640)(-0.8640,-0.5780)(-0.8232,-0.5897)(-0.7792,-0.5990)(-0.7319,-0.6058)(-0.6818,-0.6103)(-0.6289,-0.6123)(-0.5734,-0.6118)(-0.5157,-0.6088)(-0.4559,-0.6034)(-0.3942,-0.5956)(-0.3310,-0.5854)(-0.2664,-0.5728)(-0.2007,-0.5579)(-0.1343,-0.5407)(-0.0673,-0.5214)(0.0000,-0.5000)};
\draw[->,thick](axis cs:1.0,0.0)--(axis cs:1.2263,0.2263);
\draw[->,thick](axis cs:0.0,1.0)--(axis cs:-0.3036,0.8988);
\draw[->,thick](axis cs:-1.0,0.0)--(axis cs:-1.2263,-0.2263);
\draw[->,thick](axis cs:0.0,-1.0)--(axis cs:0.3036,-0.8988);
\end{axis}
\begin{axis}[at={(4.35cm,0)},width=4.3cm,height=4cm,xmin=-2,xmax=2,ymin=-2,ymax=2,axis lines=middle,xlabel={$x$},ylabel={$y$},title={$a=2$},ticks=none,clip=true]
\addplot[main,thin] coordinates{(-0.2037,-0.2668)(-0.1892,-0.2643)(-0.1741,-0.2614)(-0.1581,-0.2580)(-0.1413,-0.2542)(-0.1238,-0.2500)(-0.1055,-0.2452)(-0.0863,-0.2400)(-0.0664,-0.2342)(-0.0456,-0.2280)(-0.0240,-0.2212)(-0.0016,-0.2139)(0.0216,-0.2061)(0.0456,-0.1976)(0.0705,-0.1887)(0.0961,-0.1791)(0.1226,-0.1690)(0.1498,-0.1583)(0.1779,-0.1470)(0.2067,-0.1351)(0.2364,-0.1226)(0.2668,-0.1095)(0.2979,-0.0957)(0.3298,-0.0813)(0.3625,-0.0663)(0.3958,-0.0507)(0.4299,-0.0344)(0.4646,-0.0175)(0.5000,0.0000)(0.5360,0.0182)(0.5727,0.0370)(0.6099,0.0564)(0.6477,0.0765)(0.6860,0.0972)(0.7249,0.1186)(0.7642,0.1405)(0.8039,0.1631)(0.8440,0.1863)(0.8845,0.2101)(0.9254,0.2345)(0.9665,0.2594)(1.0078,0.2850)(1.0494,0.3111)(1.0911,0.3377)(1.1328,0.3649)(1.1747,0.3926)(1.2165,0.4207)(1.2583,0.4494)(1.3000,0.4785)(1.3414,0.5081)(1.3827,0.5380)(1.4237,0.5684)(1.4643,0.5991)(1.5044,0.6302)(1.5441,0.6615)(1.5832,0.6932)(1.6217,0.7251)(1.6595,0.7572)(1.6965,0.7895)(1.7327,0.8220)(1.7679,0.8546)(1.8021,0.8872)(1.8352,0.9199)(1.8671,0.9526)(1.8977,0.9852)(1.9270,1.0178)(1.9549,1.0502)(1.9812,1.0824)(2.0060,1.1143)(2.0290,1.1460)(2.0502,1.1774)(2.0696,1.2084)(2.0869,1.2389)(2.1022,1.2689)(2.1153,1.2984)(2.1261,1.3272)(2.1346,1.3554)(2.1406,1.3829)(2.1440,1.4095)(2.1448,1.4353)(2.1428,1.4602)(2.1379,1.4841)(2.1302,1.5069)(2.1193,1.5286)(2.1053,1.5490)};
\addplot[main,thin] coordinates{(0.8003,0.5966)(0.7929,0.6036)(0.7842,0.6101)(0.7741,0.6160)(0.7627,0.6214)(0.7499,0.6261)(0.7357,0.6302)(0.7200,0.6337)(0.7027,0.6364)(0.6840,0.6384)(0.6637,0.6396)(0.6417,0.6401)(0.6182,0.6397)(0.5929,0.6385)(0.5660,0.6365)(0.5374,0.6335)(0.5070,0.6296)(0.4749,0.6247)(0.4410,0.6189)(0.4053,0.6120)(0.3677,0.6041)(0.3284,0.5951)(0.2871,0.5850)(0.2440,0.5738)(0.1990,0.5615)(0.1521,0.5479)(0.1033,0.5332)(0.0526,0.5172)(0.0000,0.5000)(-0.0545,0.4815)(-0.1110,0.4617)(-0.1693,0.4406)(-0.2296,0.4181)(-0.2917,0.3943)(-0.3557,0.3691)(-0.4216,0.3425)(-0.4894,0.3145)(-0.5589,0.2851)(-0.6303,0.2542)(-0.7034,0.2219)(-0.7783,0.1881)(-0.8549,0.1529)(-0.9332,0.1162)(-1.0131,0.0779)(-1.0946,0.0382)(-1.1777,-0.0030)(-1.2622,-0.0457)(-1.3482,-0.0899)(-1.4355,-0.1356)(-1.5242,-0.1827)(-1.6141,-0.2314)(-1.7052,-0.2815)(-1.7973,-0.3331)(-1.8905,-0.3861)(-1.9846,-0.4405)(-2.0796,-0.4964)(-2.1753,-0.5536)(-2.2717,-0.6122)(-2.3686,-0.6721)(-2.4660,-0.7334)};
\addplot[main,thin] coordinates{(0.2037,0.2668)(0.1892,0.2643)(0.1741,0.2614)(0.1581,0.2580)(0.1413,0.2542)(0.1238,0.2500)(0.1055,0.2452)(0.0863,0.2400)(0.0664,0.2342)(0.0456,0.2280)(0.0240,0.2212)(0.0016,0.2139)(-0.0216,0.2061)(-0.0456,0.1976)(-0.0705,0.1887)(-0.0961,0.1791)(-0.1226,0.1690)(-0.1498,0.1583)(-0.1779,0.1470)(-0.2067,0.1351)(-0.2364,0.1226)(-0.2668,0.1095)(-0.2979,0.0957)(-0.3298,0.0813)(-0.3625,0.0663)(-0.3958,0.0507)(-0.4299,0.0344)(-0.4646,0.0175)(-0.5000,0.0000)(-0.5360,-0.0182)(-0.5727,-0.0370)(-0.6099,-0.0564)(-0.6477,-0.0765)(-0.6860,-0.0972)(-0.7249,-0.1186)(-0.7642,-0.1405)(-0.8039,-0.1631)(-0.8440,-0.1863)(-0.8845,-0.2101)(-0.9254,-0.2345)(-0.9665,-0.2594)(-1.0078,-0.2850)(-1.0494,-0.3111)(-1.0911,-0.3377)(-1.1328,-0.3649)(-1.1747,-0.3926)(-1.2165,-0.4207)(-1.2583,-0.4494)(-1.3000,-0.4785)(-1.3414,-0.5081)(-1.3827,-0.5380)(-1.4237,-0.5684)(-1.4643,-0.5991)(-1.5044,-0.6302)(-1.5441,-0.6615)(-1.5832,-0.6932)(-1.6217,-0.7251)(-1.6595,-0.7572)(-1.6965,-0.7895)(-1.7327,-0.8220)(-1.7679,-0.8546)(-1.8021,-0.8872)(-1.8352,-0.9199)(-1.8671,-0.9526)(-1.8977,-0.9852)(-1.9270,-1.0178)(-1.9549,-1.0502)(-1.9812,-1.0824)(-2.0060,-1.1143)(-2.0290,-1.1460)(-2.0502,-1.1774)(-2.0696,-1.2084)(-2.0869,-1.2389)(-2.1022,-1.2689)(-2.1153,-1.2984)(-2.1261,-1.3272)(-2.1346,-1.3554)(-2.1406,-1.3829)(-2.1440,-1.4095)(-2.1448,-1.4353)(-2.1428,-1.4602)(-2.1379,-1.4841)(-2.1302,-1.5069)(-2.1193,-1.5286)(-2.1053,-1.5490)};
\addplot[main,thin] coordinates{(-0.8003,-0.5966)(-0.7929,-0.6036)(-0.7842,-0.6101)(-0.7741,-0.6160)(-0.7627,-0.6214)(-0.7499,-0.6261)(-0.7357,-0.6302)(-0.7200,-0.6337)(-0.7027,-0.6364)(-0.6840,-0.6384)(-0.6637,-0.6396)(-0.6417,-0.6401)(-0.6182,-0.6397)(-0.5929,-0.6385)(-0.5660,-0.6365)(-0.5374,-0.6335)(-0.5070,-0.6296)(-0.4749,-0.6247)(-0.4410,-0.6189)(-0.4053,-0.6120)(-0.3677,-0.6041)(-0.3284,-0.5951)(-0.2871,-0.5850)(-0.2440,-0.5738)(-0.1990,-0.5615)(-0.1521,-0.5479)(-0.1033,-0.5332)(-0.0526,-0.5172)(0.0000,-0.5000)(0.0545,-0.4815)(0.1110,-0.4617)(0.1693,-0.4406)(0.2296,-0.4181)(0.2917,-0.3943)(0.3557,-0.3691)(0.4216,-0.3425)(0.4894,-0.3145)(0.5589,-0.2851)(0.6303,-0.2542)(0.7034,-0.2219)(0.7783,-0.1881)(0.8549,-0.1529)(0.9332,-0.1162)(1.0131,-0.0779)(1.0946,-0.0382)(1.1777,0.0030)(1.2622,0.0457)(1.3482,0.0899)(1.4355,0.1356)(1.5242,0.1827)(1.6141,0.2314)(1.7052,0.2815)(1.7973,0.3331)(1.8905,0.3861)(1.9846,0.4405)(2.0796,0.4964)(2.1753,0.5536)(2.2717,0.6122)(2.3686,0.6721)(2.4660,0.7334)};
\draw[->,thick](axis cs:1.0,0.0)--(axis cs:1.2862,0.1431);
\draw[->,thick](axis cs:0.0,1.0)--(axis cs:-0.3036,0.8988);
\draw[->,thick](axis cs:-1.0,0.0)--(axis cs:-1.2862,-0.1431);
\draw[->,thick](axis cs:0.0,-1.0)--(axis cs:0.3036,-0.8988);
\end{axis}
\begin{axis}[at={(8.7cm,0)},width=4.3cm,height=4cm,xmin=-2,xmax=2,ymin=-2,ymax=2,axis lines=middle,xlabel={$x$},ylabel={$y$},title={$a=a_+$},ticks=none,clip=true]
\addplot[dashed] coordinates{(2.2000,1.2702)(-2.2000,-1.2702)};
\addplot[dashed] coordinates{(2.2000,1.2702)(-2.2000,-1.2702)};
\addplot[main,thin] coordinates{(-0.1760,-0.2405)(-0.1654,-0.2380)(-0.1541,-0.2353)(-0.1421,-0.2322)(-0.1294,-0.2288)(-0.1159,-0.2251)(-0.1015,-0.2210)(-0.0863,-0.2166)(-0.0703,-0.2117)(-0.0533,-0.2065)(-0.0354,-0.2008)(-0.0165,-0.1946)(0.0034,-0.1880)(0.0244,-0.1810)(0.0465,-0.1734)(0.0697,-0.1653)(0.0941,-0.1566)(0.1198,-0.1473)(0.1468,-0.1375)(0.1752,-0.1270)(0.2049,-0.1159)(0.2361,-0.1041)(0.2688,-0.0916)(0.3030,-0.0783)(0.3389,-0.0643)(0.3765,-0.0495)(0.4158,-0.0339)(0.4570,-0.0174)(0.5000,0.0000)(0.5450,0.0183)(0.5920,0.0376)(0.6412,0.0579)(0.6925,0.0793)(0.7461,0.1018)(0.8020,0.1253)(0.8604,0.1501)(0.9213,0.1761)(0.9849,0.2034)(1.0511,0.2319)(1.1202,0.2619)(1.1922,0.2933)(1.2672,0.3261)(1.3454,0.3605)(1.4268,0.3965)(1.5116,0.4341)(1.5999,0.4735)(1.6918,0.5146)(1.7874,0.5576)(1.8870,0.6025)(1.9905,0.6493)(2.0982,0.6983)(2.2102,0.7494)(2.3267,0.8027)(2.4478,0.8583)};
\addplot[main,thin] coordinates{(0.7214,0.6570)(0.7140,0.6591)(0.7058,0.6609)(0.6966,0.6623)(0.6865,0.6633)(0.6753,0.6639)(0.6631,0.6641)(0.6497,0.6639)(0.6351,0.6631)(0.6194,0.6618)(0.6023,0.6601)(0.5839,0.6577)(0.5641,0.6548)(0.5429,0.6512)(0.5201,0.6470)(0.4958,0.6422)(0.4697,0.6366)(0.4420,0.6302)(0.4125,0.6231)(0.3811,0.6152)(0.3477,0.6064)(0.3123,0.5967)(0.2748,0.5860)(0.2350,0.5744)(0.1930,0.5618)(0.1486,0.5481)(0.1017,0.5332)(0.0522,0.5172)(0.0000,0.5000)(-0.0550,0.4815)(-0.1129,0.4617)(-0.1738,0.4404)(-0.2379,0.4178)(-0.3053,0.3936)(-0.3760,0.3678)(-0.4503,0.3404)(-0.5283,0.3113)(-0.6101,0.2804)(-0.6958,0.2477)(-0.7856,0.2130)(-0.8798,0.1763)(-0.9783,0.1376)(-1.0815,0.0966)(-1.1894,0.0534)(-1.3023,0.0078)(-1.4204,-0.0402)(-1.5438,-0.0908)(-1.6727,-0.1441)(-1.8074,-0.2000)(-1.9480,-0.2589)(-2.0949,-0.3207)(-2.2481,-0.3857)(-2.4080,-0.4538)};
\addplot[main,thin] coordinates{(0.1760,0.2405)(0.1654,0.2380)(0.1541,0.2353)(0.1421,0.2322)(0.1294,0.2288)(0.1159,0.2251)(0.1015,0.2210)(0.0863,0.2166)(0.0703,0.2117)(0.0533,0.2065)(0.0354,0.2008)(0.0165,0.1946)(-0.0034,0.1880)(-0.0244,0.1810)(-0.0465,0.1734)(-0.0697,0.1653)(-0.0941,0.1566)(-0.1198,0.1473)(-0.1468,0.1375)(-0.1752,0.1270)(-0.2049,0.1159)(-0.2361,0.1041)(-0.2688,0.0916)(-0.3030,0.0783)(-0.3389,0.0643)(-0.3765,0.0495)(-0.4158,0.0339)(-0.4570,0.0174)(-0.5000,0.0000)(-0.5450,-0.0183)(-0.5920,-0.0376)(-0.6412,-0.0579)(-0.6925,-0.0793)(-0.7461,-0.1018)(-0.8020,-0.1253)(-0.8604,-0.1501)(-0.9213,-0.1761)(-0.9849,-0.2034)(-1.0511,-0.2319)(-1.1202,-0.2619)(-1.1922,-0.2933)(-1.2672,-0.3261)(-1.3454,-0.3605)(-1.4268,-0.3965)(-1.5116,-0.4341)(-1.5999,-0.4735)(-1.6918,-0.5146)(-1.7874,-0.5576)(-1.8870,-0.6025)(-1.9905,-0.6493)(-2.0982,-0.6983)(-2.2102,-0.7494)(-2.3267,-0.8027)(-2.4478,-0.8583)};
\addplot[main,thin] coordinates{(-0.7214,-0.6570)(-0.7140,-0.6591)(-0.7058,-0.6609)(-0.6966,-0.6623)(-0.6865,-0.6633)(-0.6753,-0.6639)(-0.6631,-0.6641)(-0.6497,-0.6639)(-0.6351,-0.6631)(-0.6194,-0.6618)(-0.6023,-0.6601)(-0.5839,-0.6577)(-0.5641,-0.6548)(-0.5429,-0.6512)(-0.5201,-0.6470)(-0.4958,-0.6422)(-0.4697,-0.6366)(-0.4420,-0.6302)(-0.4125,-0.6231)(-0.3811,-0.6152)(-0.3477,-0.6064)(-0.3123,-0.5967)(-0.2748,-0.5860)(-0.2350,-0.5744)(-0.1930,-0.5618)(-0.1486,-0.5481)(-0.1017,-0.5332)(-0.0522,-0.5172)(0.0000,-0.5000)(0.0550,-0.4815)(0.1129,-0.4617)(0.1738,-0.4404)(0.2379,-0.4178)(0.3053,-0.3936)(0.3760,-0.3678)(0.4503,-0.3404)(0.5283,-0.3113)(0.6101,-0.2804)(0.6958,-0.2477)(0.7856,-0.2130)(0.8798,-0.1763)(0.9783,-0.1376)(1.0815,-0.0966)(1.1894,-0.0534)(1.3023,-0.0078)(1.4204,0.0402)(1.5438,0.0908)(1.6727,0.1441)(1.8074,0.2000)(1.9480,0.2589)(2.0949,0.3207)(2.2481,0.3857)(2.4080,0.4538)};
\draw[->,thick](axis cs:1.0,0.0)--(axis cs:1.2965,0.1203);
\draw[->,thick](axis cs:0.0,1.0)--(axis cs:-0.3036,0.8988);
\draw[->,thick](axis cs:-1.0,0.0)--(axis cs:-1.2965,-0.1203);
\draw[->,thick](axis cs:0.0,-1.0)--(axis cs:0.3036,-0.8988);
\end{axis}
\end{tikzpicture}
\par
\begin{tikzpicture}
\begin{axis}[at={(0.0cm,0)},width=4.3cm,height=4cm,xmin=-2,xmax=2,ymin=-2,ymax=2,axis lines=middle,xlabel={$x$},ylabel={$y$},title={$a=11/4$},ticks=none,clip=true]
\addplot[dashed] coordinates{(2.2000,0.8484)(-2.2000,-0.8484)};
\addplot[dashed] coordinates{(2.2000,1.9016)(-2.2000,-1.9016)};
\addplot[main,thin] coordinates{(-0.1608,-0.2268)(-0.1519,-0.2243)(-0.1425,-0.2216)(-0.1324,-0.2187)(-0.1217,-0.2155)(-0.1101,-0.2120)(-0.0979,-0.2082)(-0.0847,-0.2041)(-0.0707,-0.1997)(-0.0558,-0.1949)(-0.0399,-0.1897)(-0.0230,-0.1842)(-0.0050,-0.1782)(0.0142,-0.1718)(0.0347,-0.1649)(0.0564,-0.1575)(0.0795,-0.1496)(0.1041,-0.1411)(0.1302,-0.1321)(0.1580,-0.1224)(0.1876,-0.1120)(0.2189,-0.1010)(0.2522,-0.0892)(0.2876,-0.0766)(0.3252,-0.0631)(0.3651,-0.0488)(0.4074,-0.0336)(0.4523,-0.0173)(0.5000,0.0000)(0.5506,0.0184)(0.6043,0.0380)(0.6612,0.0589)(0.7216,0.0811)(0.7856,0.1047)(0.8535,0.1297)(0.9256,0.1564)(1.0019,0.1847)(1.0829,0.2148)(1.1687,0.2468)(1.2597,0.2807)(1.3561,0.3167)(1.4583,0.3550)(1.5666,0.3956)(1.6814,0.4387)(1.8030,0.4844)(1.9318,0.5329)(2.0684,0.5844)(2.2131,0.6390)(2.3663,0.6969)};
\addplot[main,thin] coordinates{(0.6804,0.6898)(0.6730,0.6893)(0.6649,0.6886)(0.6560,0.6876)(0.6464,0.6864)(0.6359,0.6848)(0.6246,0.6829)(0.6123,0.6806)(0.5990,0.6780)(0.5846,0.6749)(0.5692,0.6715)(0.5525,0.6676)(0.5346,0.6633)(0.5154,0.6584)(0.4947,0.6530)(0.4726,0.6471)(0.4488,0.6406)(0.4234,0.6334)(0.3962,0.6256)(0.3672,0.6170)(0.3361,0.6077)(0.3029,0.5976)(0.2675,0.5866)(0.2297,0.5748)(0.1894,0.5620)(0.1465,0.5482)(0.1007,0.5333)(0.0519,0.5172)(0.0000,0.5000)(-0.0553,0.4815)(-0.1141,0.4616)(-0.1767,0.4403)(-0.2432,0.4175)(-0.3140,0.3931)(-0.3892,0.3670)(-0.4692,0.3391)(-0.5542,0.3092)(-0.6444,0.2774)(-0.7403,0.2433)(-0.8421,0.2070)(-0.9502,0.1684)(-1.0649,0.1271)(-1.1867,0.0832)(-1.3160,0.0363)(-1.4532,-0.0135)(-1.5987,-0.0666)(-1.7531,-0.1230)(-1.9169,-0.1831)(-2.0906,-0.2469)(-2.2748,-0.3148)(-2.4702,-0.3870)};
\addplot[main,thin] coordinates{(0.1608,0.2268)(0.1519,0.2243)(0.1425,0.2216)(0.1324,0.2187)(0.1217,0.2155)(0.1101,0.2120)(0.0979,0.2082)(0.0847,0.2041)(0.0707,0.1997)(0.0558,0.1949)(0.0399,0.1897)(0.0230,0.1842)(0.0050,0.1782)(-0.0142,0.1718)(-0.0347,0.1649)(-0.0564,0.1575)(-0.0795,0.1496)(-0.1041,0.1411)(-0.1302,0.1321)(-0.1580,0.1224)(-0.1876,0.1120)(-0.2189,0.1010)(-0.2522,0.0892)(-0.2876,0.0766)(-0.3252,0.0631)(-0.3651,0.0488)(-0.4074,0.0336)(-0.4523,0.0173)(-0.5000,0.0000)(-0.5506,-0.0184)(-0.6043,-0.0380)(-0.6612,-0.0589)(-0.7216,-0.0811)(-0.7856,-0.1047)(-0.8535,-0.1297)(-0.9256,-0.1564)(-1.0019,-0.1847)(-1.0829,-0.2148)(-1.1687,-0.2468)(-1.2597,-0.2807)(-1.3561,-0.3167)(-1.4583,-0.3550)(-1.5666,-0.3956)(-1.6814,-0.4387)(-1.8030,-0.4844)(-1.9318,-0.5329)(-2.0684,-0.5844)(-2.2131,-0.6390)(-2.3663,-0.6969)};
\addplot[main,thin] coordinates{(-0.6804,-0.6898)(-0.6730,-0.6893)(-0.6649,-0.6886)(-0.6560,-0.6876)(-0.6464,-0.6864)(-0.6359,-0.6848)(-0.6246,-0.6829)(-0.6123,-0.6806)(-0.5990,-0.6780)(-0.5846,-0.6749)(-0.5692,-0.6715)(-0.5525,-0.6676)(-0.5346,-0.6633)(-0.5154,-0.6584)(-0.4947,-0.6530)(-0.4726,-0.6471)(-0.4488,-0.6406)(-0.4234,-0.6334)(-0.3962,-0.6256)(-0.3672,-0.6170)(-0.3361,-0.6077)(-0.3029,-0.5976)(-0.2675,-0.5866)(-0.2297,-0.5748)(-0.1894,-0.5620)(-0.1465,-0.5482)(-0.1007,-0.5333)(-0.0519,-0.5172)(0.0000,-0.5000)(0.0553,-0.4815)(0.1141,-0.4616)(0.1767,-0.4403)(0.2432,-0.4175)(0.3140,-0.3931)(0.3892,-0.3670)(0.4692,-0.3391)(0.5542,-0.3092)(0.6444,-0.2774)(0.7403,-0.2433)(0.8421,-0.2070)(0.9502,-0.1684)(1.0649,-0.1271)(1.1867,-0.0832)(1.3160,-0.0363)(1.4532,0.0135)(1.5987,0.0666)(1.7531,0.1230)(1.9169,0.1831)(2.0906,0.2469)(2.2748,0.3148)(2.4702,0.3870)};
\draw[->,thick](axis cs:1.0,0.0)--(axis cs:1.3007,0.1094);
\draw[->,thick](axis cs:0.0,1.0)--(axis cs:-0.3036,0.8988);
\draw[->,thick](axis cs:-1.0,0.0)--(axis cs:-1.3007,-0.1094);
\draw[->,thick](axis cs:0.0,-1.0)--(axis cs:0.3036,-0.8988);
\end{axis}
\begin{axis}[at={(4.35cm,0)},width=4.3cm,height=4cm,xmin=-2,xmax=2,ymin=-2,ymax=2,axis lines=middle,xlabel={$x$},ylabel={$y$},title={$a=3$},ticks=none,clip=true]
\addplot[dashed] coordinates{(2.2000,0.7333)(-2.2000,-0.7333)};
\addplot[dashed] coordinates{(2.2000,2.2000)(-2.2000,-2.2000)};
\addplot[main,thin] coordinates{(-0.1485,-0.2162)(-0.1410,-0.2137)(-0.1329,-0.2110)(-0.1242,-0.2081)(-0.1149,-0.2050)(-0.1049,-0.2016)(-0.0942,-0.1981)(-0.0827,-0.1942)(-0.0703,-0.1901)(-0.0570,-0.1857)(-0.0427,-0.1809)(-0.0273,-0.1758)(-0.0108,-0.1703)(0.0069,-0.1644)(0.0259,-0.1580)(0.0463,-0.1512)(0.0683,-0.1439)(0.0918,-0.1361)(0.1172,-0.1276)(0.1443,-0.1186)(0.1735,-0.1088)(0.2049,-0.0984)(0.2386,-0.0871)(0.2748,-0.0751)(0.3136,-0.0621)(0.3553,-0.0482)(0.4002,-0.0333)(0.4483,-0.0172)(0.5000,0.0000)(0.5555,0.0185)(0.6152,0.0384)(0.6792,0.0597)(0.7480,0.0827)(0.8219,0.1073)(0.9013,0.1338)(0.9865,0.1622)(1.0781,0.1927)(1.1764,0.2255)(1.2820,0.2607)(1.3955,0.2985)(1.5173,0.3391)(1.6482,0.3827)(1.7887,0.4296)(1.9397,0.4799)(2.1018,0.5339)(2.2759,0.5920)(2.4629,0.6543)};
\addplot[main,thin] coordinates{(0.6485,0.7162)(0.6410,0.7137)(0.6329,0.7110)(0.6242,0.7081)(0.6149,0.7050)(0.6049,0.7016)(0.5942,0.6981)(0.5827,0.6942)(0.5703,0.6901)(0.5570,0.6857)(0.5427,0.6809)(0.5273,0.6758)(0.5108,0.6703)(0.4931,0.6644)(0.4741,0.6580)(0.4537,0.6512)(0.4317,0.6439)(0.4082,0.6361)(0.3828,0.6276)(0.3557,0.6186)(0.3265,0.6088)(0.2951,0.5984)(0.2614,0.5871)(0.2252,0.5751)(0.1864,0.5621)(0.1447,0.5482)(0.0998,0.5333)(0.0517,0.5172)(0.0000,0.5000)(-0.0555,0.4815)(-0.1152,0.4616)(-0.1792,0.4403)(-0.2480,0.4173)(-0.3219,0.3927)(-0.4013,0.3662)(-0.4865,0.3378)(-0.5781,0.3073)(-0.6764,0.2745)(-0.7820,0.2393)(-0.8955,0.2015)(-1.0173,0.1609)(-1.1482,0.1173)(-1.2887,0.0704)(-1.4397,0.0201)(-1.6018,-0.0339)(-1.7759,-0.0920)(-1.9629,-0.1543)(-2.1638,-0.2213)(-2.3796,-0.2932)};
\addplot[main,thin] coordinates{(0.1485,0.2162)(0.1410,0.2137)(0.1329,0.2110)(0.1242,0.2081)(0.1149,0.2050)(0.1049,0.2016)(0.0942,0.1981)(0.0827,0.1942)(0.0703,0.1901)(0.0570,0.1857)(0.0427,0.1809)(0.0273,0.1758)(0.0108,0.1703)(-0.0069,0.1644)(-0.0259,0.1580)(-0.0463,0.1512)(-0.0683,0.1439)(-0.0918,0.1361)(-0.1172,0.1276)(-0.1443,0.1186)(-0.1735,0.1088)(-0.2049,0.0984)(-0.2386,0.0871)(-0.2748,0.0751)(-0.3136,0.0621)(-0.3553,0.0482)(-0.4002,0.0333)(-0.4483,0.0172)(-0.5000,0.0000)(-0.5555,-0.0185)(-0.6152,-0.0384)(-0.6792,-0.0597)(-0.7480,-0.0827)(-0.8219,-0.1073)(-0.9013,-0.1338)(-0.9865,-0.1622)(-1.0781,-0.1927)(-1.1764,-0.2255)(-1.2820,-0.2607)(-1.3955,-0.2985)(-1.5173,-0.3391)(-1.6482,-0.3827)(-1.7887,-0.4296)(-1.9397,-0.4799)(-2.1018,-0.5339)(-2.2759,-0.5920)(-2.4629,-0.6543)};
\addplot[main,thin] coordinates{(-0.6485,-0.7162)(-0.6410,-0.7137)(-0.6329,-0.7110)(-0.6242,-0.7081)(-0.6149,-0.7050)(-0.6049,-0.7016)(-0.5942,-0.6981)(-0.5827,-0.6942)(-0.5703,-0.6901)(-0.5570,-0.6857)(-0.5427,-0.6809)(-0.5273,-0.6758)(-0.5108,-0.6703)(-0.4931,-0.6644)(-0.4741,-0.6580)(-0.4537,-0.6512)(-0.4317,-0.6439)(-0.4082,-0.6361)(-0.3828,-0.6276)(-0.3557,-0.6186)(-0.3265,-0.6088)(-0.2951,-0.5984)(-0.2614,-0.5871)(-0.2252,-0.5751)(-0.1864,-0.5621)(-0.1447,-0.5482)(-0.0998,-0.5333)(-0.0517,-0.5172)(0.0000,-0.5000)(0.0555,-0.4815)(0.1152,-0.4616)(0.1792,-0.4403)(0.2480,-0.4173)(0.3219,-0.3927)(0.4013,-0.3662)(0.4865,-0.3378)(0.5781,-0.3073)(0.6764,-0.2745)(0.7820,-0.2393)(0.8955,-0.2015)(1.0173,-0.1609)(1.1482,-0.1173)(1.2887,-0.0704)(1.4397,-0.0201)(1.6018,0.0339)(1.7759,0.0920)(1.9629,0.1543)(2.1638,0.2213)(2.3796,0.2932)};
\draw[->,thick](axis cs:1.0,0.0)--(axis cs:1.3036,0.1012);
\draw[->,thick](axis cs:0.0,1.0)--(axis cs:-0.3036,0.8988);
\draw[->,thick](axis cs:-1.0,0.0)--(axis cs:-1.3036,-0.1012);
\draw[->,thick](axis cs:0.0,-1.0)--(axis cs:0.3036,-0.8988);
\end{axis}
\begin{axis}[at={(8.7cm,0)},width=4.3cm,height=4cm,xmin=-2,xmax=2,ymin=-2,ymax=2,axis lines=middle,xlabel={$x$},ylabel={$y$},title={$a=4$},ticks=none,clip=true]
\addplot[dashed] coordinates{(2.2000,0.5113)(-2.2000,-0.5113)};
\addplot[dashed] coordinates{(1.5339,2.2000)(-1.5339,-2.2000)};
\addplot[main,thin] coordinates{(-0.1089,-0.1826)(-0.1048,-0.1800)(-0.1003,-0.1772)(-0.0954,-0.1745)(-0.0902,-0.1716)(-0.0844,-0.1686)(-0.0781,-0.1656)(-0.0712,-0.1624)(-0.0636,-0.1591)(-0.0553,-0.1556)(-0.0461,-0.1519)(-0.0359,-0.1480)(-0.0246,-0.1439)(-0.0120,-0.1395)(0.0019,-0.1347)(0.0175,-0.1297)(0.0348,-0.1242)(0.0541,-0.1183)(0.0757,-0.1119)(0.0998,-0.1049)(0.1268,-0.0972)(0.1570,-0.0888)(0.1909,-0.0796)(0.2288,-0.0695)(0.2713,-0.0583)(0.3190,-0.0459)(0.3725,-0.0322)(0.4326,-0.0169)(0.5000,0.0000)(0.5757,0.0189)(0.6608,0.0399)(0.7564,0.0633)(0.8638,0.0895)(0.9846,0.1187)(1.1203,0.1515)(1.2729,0.1881)(1.4444,0.2291)(1.6373,0.2751)(1.8542,0.3266)(2.0982,0.3845)(2.3725,0.4493)};
\addplot[main,thin] coordinates{(0.5478,0.8041)(0.5399,0.7950)(0.5317,0.7859)(0.5234,0.7768)(0.5148,0.7678)(0.5059,0.7588)(0.4967,0.7497)(0.4871,0.7407)(0.4772,0.7316)(0.4667,0.7225)(0.4556,0.7133)(0.4440,0.7041)(0.4316,0.6947)(0.4184,0.6853)(0.4042,0.6757)(0.3890,0.6659)(0.3727,0.6559)(0.3549,0.6456)(0.3356,0.6351)(0.3146,0.6242)(0.2917,0.6130)(0.2665,0.6013)(0.2389,0.5890)(0.2085,0.5762)(0.1749,0.5627)(0.1377,0.5485)(0.0965,0.5334)(0.0508,0.5172)(0.0000,0.5000)(-0.0566,0.4815)(-0.1196,0.4615)(-0.1899,0.4399)(-0.2684,0.4164)(-0.3562,0.3909)(-0.4544,0.3630)(-0.5643,0.3324)(-0.6874,0.2988)(-0.8253,0.2618)(-0.9799,0.2210)(-1.1534,0.1759)(-1.3480,0.1259)(-1.5664,0.0704)(-1.8117,0.0087)(-2.0871,-0.0599)(-2.3965,-0.1364)};
\addplot[main,thin] coordinates{(0.1089,0.1826)(0.1048,0.1800)(0.1003,0.1772)(0.0954,0.1745)(0.0902,0.1716)(0.0844,0.1686)(0.0781,0.1656)(0.0712,0.1624)(0.0636,0.1591)(0.0553,0.1556)(0.0461,0.1519)(0.0359,0.1480)(0.0246,0.1439)(0.0120,0.1395)(-0.0019,0.1347)(-0.0175,0.1297)(-0.0348,0.1242)(-0.0541,0.1183)(-0.0757,0.1119)(-0.0998,0.1049)(-0.1268,0.0972)(-0.1570,0.0888)(-0.1909,0.0796)(-0.2288,0.0695)(-0.2713,0.0583)(-0.3190,0.0459)(-0.3725,0.0322)(-0.4326,0.0169)(-0.5000,0.0000)(-0.5757,-0.0189)(-0.6608,-0.0399)(-0.7564,-0.0633)(-0.8638,-0.0895)(-0.9846,-0.1187)(-1.1203,-0.1515)(-1.2729,-0.1881)(-1.4444,-0.2291)(-1.6373,-0.2751)(-1.8542,-0.3266)(-2.0982,-0.3845)(-2.3725,-0.4493)};
\addplot[main,thin] coordinates{(-0.5478,-0.8041)(-0.5399,-0.7950)(-0.5317,-0.7859)(-0.5234,-0.7768)(-0.5148,-0.7678)(-0.5059,-0.7588)(-0.4967,-0.7497)(-0.4871,-0.7407)(-0.4772,-0.7316)(-0.4667,-0.7225)(-0.4556,-0.7133)(-0.4440,-0.7041)(-0.4316,-0.6947)(-0.4184,-0.6853)(-0.4042,-0.6757)(-0.3890,-0.6659)(-0.3727,-0.6559)(-0.3549,-0.6456)(-0.3356,-0.6351)(-0.3146,-0.6242)(-0.2917,-0.6130)(-0.2665,-0.6013)(-0.2389,-0.5890)(-0.2085,-0.5762)(-0.1749,-0.5627)(-0.1377,-0.5485)(-0.0965,-0.5334)(-0.0508,-0.5172)(0.0000,-0.5000)(0.0566,-0.4815)(0.1196,-0.4615)(0.1899,-0.4399)(0.2684,-0.4164)(0.3562,-0.3909)(0.4544,-0.3630)(0.5643,-0.3324)(0.6874,-0.2988)(0.8253,-0.2618)(0.9799,-0.2210)(1.1534,-0.1759)(1.3480,-0.1259)(1.5664,-0.0704)(1.8117,-0.0087)(2.0871,0.0599)(2.3965,0.1364)};
\draw[->,thick](axis cs:1.0,0.0)--(axis cs:1.3104,0.0776);
\draw[->,thick](axis cs:0.0,1.0)--(axis cs:-0.3036,0.8988);
\draw[->,thick](axis cs:-1.0,0.0)--(axis cs:-1.3104,-0.0776);
\draw[->,thick](axis cs:0.0,-1.0)--(axis cs:0.3036,-0.8988);
\end{axis}
\end{tikzpicture}
\par
\end{center}
\end{solution}
\begin{exercise}\label{cw:ps9-II-36}原题 II.36（矩阵指数）：
(a) 如习题课 21 第 4 题，复数 $z=a+bi$ 对应 $A(z)=\begin{pmatrix}a&-b\\b&a\end{pmatrix}$，它表示乘法：$z(x+yi)=v+wi$ 时 $A(z)(x,y)^T=(v,w)^T$。求 $e^{A(z)t}$、$A(e^{zt})$。
(b) 若 $mx''+bx'+kx=0$ 的一对解满足 $x_1(0)=1,x_1'(0)=0,x_2(0)=0,x_2'(0)=1$，称在 $0$ 归一化。求 $x''+2x'+2x=0$ 的归一化解对，再求算子 $D^2+2D+2I$ 伴随矩阵的 $e^{At}$，作评论。
(c) 已知 $e^{3t}(1,1)^T$ 与 $e^{2t}(1,2)^T$ 满足 $\mathbf u'=A\mathbf u$：(i) 求 $\mathbf u_1(0)=(1,0)^T,\mathbf u_2(0)=(0,1)^T$ 的两个解；(ii) 求 $e^{At}$；(iii) 求 $A$。\end{exercise}
\begin{solution}
\textit{官方解译审与补证（AI 辅助）；其中另标编者补解或原文修正。}\par

(a) 记 $J=\begin{pmatrix}0&-1\\1&0\end{pmatrix}$，则 $J^2=-I$、$A=aI+bJ$。两部分可交换，故
\[e^{A(z)t}=e^{at}(I\cos bt+J\sin bt)=e^{at}\begin{pmatrix}\cos bt&-\sin bt\\\sin bt&\cos bt\end{pmatrix}.\]
另一方面 $e^{zt}=e^{at}(\cos bt+i\sin bt)$，代入映射 $A$ 得到同一矩阵。因此 $A(e^{zt})=e^{A(z)t}$，包括 $b=0$ 的情况。此等式反映复乘法与其二维实矩阵表示相容。
(b) 基本解 $y_1=e^{-t}\cos t,y_2=e^{-t}\sin t$ 的初值为 $(1,-1)$、$(0,1)$，所以
\[x_1=y_1+y_2=e^{-t}(\cos t+\sin t),\qquad x_2=y_2=e^{-t}\sin t.\]
$A=\begin{pmatrix}0&1\\-2&-2\end{pmatrix}$，将 $x_j$ 及 $x_j'$ 作第 $j$ 列，得
\[e^{At}=\begin{pmatrix}x_1&x_2\\x_1'&x_2'\end{pmatrix}=e^{-t}\begin{pmatrix}\cos t+\sin t&\sin t\\-2\sin t&\cos t-\sin t\end{pmatrix}.\]
此矩阵在 $t=0$ 为 $I$，各列解伴随系统，故由初值唯一性确为矩阵指数。一般基本解矩阵若不在 $0$ 等于 $I$，不能直接当 $e^{At}$；归一化正是将两列分别对应于两个坐标初值。
(c)(i) 令 $\mathbf u=c_1e^{3t}(1,1)^T+c_2e^{2t}(1,2)^T$，零点值为 $(c_1+c_2,c_1+2c_2)^T$。第一个初值得 $c_1=2,c_2=-1$，第二个得 $c_1=-1,c_2=1$，于是
\[\mathbf u_1=\binom{2e^{3t}-e^{2t}}{2e^{3t}-2e^{2t}},\qquad\mathbf u_2=\binom{-e^{3t}+e^{2t}}{-e^{3t}+2e^{2t}}.\]
(ii) 两列构成在 $0$ 归一化的基本矩阵：
\[e^{At}=\begin{pmatrix}2e^{3t}-e^{2t}&-e^{3t}+e^{2t}\\2e^{3t}-2e^{2t}&-e^{3t}+2e^{2t}\end{pmatrix}.\]
(iii) 在 $0$ 求导得 $A=\begin{pmatrix}4&-1\\2&1\end{pmatrix}$。亦可令 $S=\begin{pmatrix}1&1\\1&2\end{pmatrix}$，由 $AS=S\operatorname{diag}(3,2)$ 得 $A=S\operatorname{diag}(3,2)S^{-1}$，并核实两个给定模态确实满足系统。
\end{solution}
